\documentclass[11pt]{amsart}

\usepackage{amsmath,amssymb,amsthm}
\usepackage[margin=1in]{geometry}
\usepackage[hyphens]{url}
\usepackage{hyperref}
\hypersetup{colorlinks=true,linkcolor=blue,citecolor=blue,urlcolor=blue}

\theoremstyle{plain}
\newtheorem{theorem}{Teorema}
\newtheorem{proposition}[theorem]{Proposição}
\newtheorem{lemma}[theorem]{Lema}
\newtheorem{corollary}[theorem]{Corolário}
\newtheorem*{conjecture}{Conjectura 3 de Wirsching}
\theoremstyle{definition}
\newtheorem{definition}[theorem]{Definição}
\theoremstyle{remark}
\newtheorem{remark}[theorem]{Observação}
\renewcommand{\proofname}{Demonstração}
\renewcommand{\refname}{Referências}

\DeclareMathOperator{\artanh}{artanh}
\newcommand{\R}{\mathbb{R}}
\newcommand{\C}{\mathbb{C}}
\newcommand{\Z}{\mathbb{Z}}
\newcommand{\N}{\mathbb{N}}

\title[Conjectura 3 de Wirsching]{A Conjectura 3 de Wirsching, e a correção periódica à assintótica de densidade de Berg--Kr\"uppel}

\author{Renato Augusto Tavares}

\begin{document}

\maketitle

\vspace{-2.5em}
\begin{center}
\normalsize Universidade Federal de Goi\'as \par
\texttt{rat@discente.ufg.br} \par
ORCID: \url{https://orcid.org/0009-0002-0196-3311}
\end{center}

\begin{abstract}
O argumento de Wirsching de 2003 para densidade positiva de predecessores sob o mapa $3n+1$ se
reduz a três conjecturas. A Conjectura 3 trata da assintótica de $\varphi$, uma densidade
invariante, próximo de zero, comparada a uma assintótica explícita $\varphi_0$ devida a Berg e
Kr\"uppel. A transformada log-Laplace exata de $\varphi$ se decompõe como uma parte suave, uma
correção periódica $H$, e um resto duplamente exponencialmente pequeno. Berg e Kr\"uppel deram essa
correção, para esta autofunção, como um produto infinito em 1998, cinco anos antes da conjectura de
Wirsching; eles não perguntaram se ela é constante. Nós damos $H$, em vez disso, como uma série de
Fourier com coeficientes em forma fechada em $\Gamma$ e $\zeta$: um teorema clássico de não-anulação
de $\zeta$ mostra diretamente que ela não é constante, e nós envolvemos rigorosamente sua oscilação.
A própria classe de comparação de Wirsching fixa uma única fase de $H$; ali, a Conjectura 3 vale,
com um limite explícito. Fora dessa classe a fase varre um período completo, de modo que a
assintótica mais forte e irrestrita $\varphi(t) \sim \kappa\,\varphi_0(t)$ quando $t \to 0^+$ falha.
Wirsching precisa de menos do que o enunciado literal da Conjectura 3; nós verificamos diretamente
que sua exigência real vale com ampla margem. Isso resolve apenas a Conjectura 3; as outras duas
conjecturas de Wirsching, e a própria densidade positiva de predecessores, permanecem em aberto.
\end{abstract}

\noindent\textbf{Palavras-chave:} mapa de Collatz, densidade de predecessores, densidade invariante,
assintótica log-periódica, transformada de Mellin, função atômica.

\noindent\textbf{Classificação de Assuntos Matemáticos 2020:} 11B83, 39B22, 44A15, 41A60, 11M06.

\section{Introdução}

Seja $T$ o mapa $3n+1$ nos inteiros positivos, $T(n) = n/2$ para $n$ par e $T(n) = (3n+1)/2$ para
$n$ ímpar. A monografia de Wirsching \cite{wirsching1998} desenvolve a maquinaria de conjuntos de
predecessores da qual descende o argumento deste artigo; \cite{wirsching2003} a leva adiante,
estudando a densidade de conjuntos de predecessores sob $T$ por meio de uma construção de
médias sobre os inteiros $3$-ádicos: uma família de operadores $S_l$, construída a partir de
contagens de caminhos no grafo de predecessores, converge na topologia de operadores fortes a um
operador limite $S_\infty$ cuja parte essencial é uma cadeia de Markov em $[0,1]$. Essa cadeia tem
uma única densidade invariante $\varphi$, e a prova de Wirsching de densidade positiva de
predecessores se reduz a três conjecturas sobre $\varphi$ e quantidades relacionadas, enunciadas
como três lacunas não fechadas no argumento.

Dessas três, este artigo resolve a Conjectura 3. Ela trata do comportamento de $\varphi$ quando seu
argumento tende a $0$, comparado a uma função de comparação explícita $\varphi_0$: a própria
equação~(7.11) de Wirsching define $\varphi_0$ como sendo exatamente a solução assintótica de Berg e
Kr\"uppel, nomeando \cite{bergkruppel1998} na frase que a introduz, de modo que a identificação é
escolha do próprio Wirsching, não uma substituição feita aqui. Berg e Kr\"uppel haviam produzido
$\varphi_0$ cinco anos antes, por transformada de Laplace e análise de ponto de sela, como a solução
assintótica de uma versão truncada da equação funcional que $\varphi$ satisfaz. O argumento deles
para a comparabilidade assintótica entre $\varphi$ e $\varphi_0$ era informal: eles dão a correção
periódica que relaciona as duas como um produto infinito explícito (Corolário~\ref{cor:bk} abaixo)
mas nunca o avaliam, de modo que tanto se uma única constante relaciona $\varphi$ a $\varphi_0$ ao
longo de toda sequência, quanto se a razão sequer converge ao longo de uma sequência arbitrária,
ficam em aberto. Wirsching considera diretamente a leitura mais forte possível dessa lacuna:

\begin{quote}
``Essa substituição seria possível se, por exemplo, $\varphi \sim c\,\varphi_0$ para alguma
constante $c > 0$. Acontece que a seguinte conjectura (ligeiramente mais fraca) \dots basta.''
\end{quote}

Ele então enuncia (\cite{wirsching2003}, imediatamente após a frase citada, com notação
normalizada: seu próprio $\delta_5$ é escrito $\delta$ ao longo deste artigo, já que nenhuma outra
classe indexada por $\delta$ dele aparece aqui, e sua constante $c$, mantida abaixo por fidelidade
ao seu enunciado, não tem relação com $c:=\log3$ introduzido na Seção~\ref{sec:periodic}):

\begin{conjecture}
Existem constantes reais positivas $c$ e $\delta$ tais que
\[
\lim_{l\to\infty} \frac{\varphi(z_l)}{\varphi_0(z_l)} = c > 0 \qquad \text{uniformemente para
sequências } (z_l) \in \widetilde A_\delta.
\]
\end{conjecture}

Aqui $\widetilde A_\delta$ é a classe de Wirsching de sequências que tendem a $0$ na taxa limítrofe
$z_l \sim l\cdot 3^{-l}$ (suas equações~(1.5) e~(3.2), reproduzidas antes da prova do
Teorema~\ref{thm:conj3}): uma sequência $(x_l)$ pertence a $\widetilde A_\delta$ exatamente quando
$\lfloor 3^lx_l\rfloor = k_l$ para uma sequência de inteiros $(k_l)$ com $|l-k_l|\le\delta\sqrt l$
para todo $l$. Isso é mais fraco do que a assintótica irrestrita que ele considera primeiro. A
constante se cancela na única desigualdade de que sua prova de fato precisa, de modo que um limite
dependente da classe já basta. Nenhum dos dois enunciados havia sido resolvido antes dos dois
teoremas a seguir.

\begin{theorem}[Conjectura 3]
\label{thm:conj3}
A Conjectura 3 de Wirsching (enunciada acima) é verdadeira. Para todo $\delta>0$, na sua classe de
comparação $\widetilde{A}_\delta$, a razão $\varphi(z_l)/\varphi_0(z_l)$ converge uniformemente para
$e^{H(0)}$, um número real bem definido pelo Teorema~\ref{thm:H} independentemente de sua expansão
decimal, avaliado aqui como $e^{H(0)} = 0.534122\ldots$ (uma avaliação em ponto flutuante; a prova
da Proposição~\ref{prop:Hcert} certifica rigorosamente a diferença próxima
$H(0)-H(\log\tfrac32)$, mas não $H(0)$ isoladamente), onde $H$ é a função periódica do
Teorema~\ref{thm:H} abaixo e a normalização de $\varphi_0$ é a própria de Berg e Kr\"uppel,
equação~(9.6) de \cite{bergkruppel1998}. A própria Conjectura 3 só afirma a existência de alguma
constante positiva, independentemente de normalização; o valor específico $e^{H(0)}$ está ligado a
esta normalização particular de $\varphi_0$.
\end{theorem}

\begin{theorem}[Falha da assintótica irrestrita]
\label{thm:unrestricted}
A assintótica mais forte $\varphi(t) \sim \kappa\,\varphi_0(t)$ quando $t \to 0^+$, para qualquer
constante única $\kappa$ e sem restrição a uma classe de comparação, é falsa. A razão
$\varphi(t)/\varphi_0(t)$ oscila indefinidamente quando $t \to 0^+$, com amplitude limitada
inferiormente por uma constante positiva: ela é governada pela função exatamente $c$-periódica $H$
avaliada na localização de sela $w_0(\tau)$, $\tau=\log(1/t)$, cuja fase módulo $c$ varre todo valor
infinitas vezes quando $\tau\to\infty$ (Lema~\ref{lem:saddles}). Que a amplitude seja positiva segue
de um teorema clássico, sem nenhum cálculo; seu tamanho numérico explícito é um resultado assistido
por computador, um certificado rigoroso de aritmética de intervalos (Proposição~\ref{prop:Hcert}).
\end{theorem}

Um único cálculo fundamenta os dois teoremas. Escreva $t = e^{-\tau}$ e seja $H(w)$ a função
$\log 3$-periódica que aparece na decomposição exata da transformada log-Laplace de $\varphi$
(Teorema~\ref{thm:H}); a ponte de ponto de sela das Seções~\ref{sec:formulaA}, \ref{sec:envelope} e
\ref{sec:bk} mostra $\log\varphi(t)-\log\varphi_0(t) = H(w_0(\tau))+o(1)$, onde a localização de
sela $w_0(\tau)$ avança por $c$ a cada período de $\tau$, a menos de uma correção que se anula:
$w_0(\tau+c)-w_0(\tau) = c+O(1/\tau)$ (Lema~\ref{lem:saddles}). A classe de comparação de Wirsching
restringe a fase de $w_0(\tau_l)$ módulo $\log 3$ a convergir para um único valor quando
$l\to\infty$; o intervalo irrestrito $\tau\to\infty$ sem restrição de classe, em vez disso, varre
essa fase por todo valor, já que $\Phi_0$ é uma bijeção (novamente Lema~\ref{lem:saddles}). Se a
razão converge ou oscila depende de ao longo de quais sequências se compara, e os dois teoremas
acima leem o mesmo objeto em dois escopos diferentes.

\begin{corollary}
\label{cor:bk}
$e^H$ é o próprio fator periódico de Berg e Kr\"uppel para a autofunção $\varphi$, da
Proposição~9.3 de \cite{bergkruppel1998}, que eles dão como um produto infinito; $H$ em si é seu
logaritmo. Sua série de Fourier, Proposição~\ref{prop:Hfourier} abaixo, é nova\footnote{Uma busca na
literatura por trabalhos sobre esta família de equações de diferença de duas escalas publicados
após \cite{bergkruppel1998}, incluindo artigos posteriores de seus próprios autores, não encontrou
nenhum que dê uma expansão de Fourier fechada para o fator periódico.}. Este corolário tem duas
partes de profundidade diferente: que $e^H$ coincide com o produto deles \emph{como função},
provado na Observação~\ref{rem:bkproduct} abaixo por álgebra pura ($e^H=e^{-Q}e^Le^{-\Delta}$
recombinando as próprias peças do Teorema~\ref{thm:H}), e que essa função é o fator periódico que
multiplica a própria assintótica de $\varphi$, o que precisa da cadeia completa de ponto de sela do
Teorema~\ref{thm:formulaA} e das Proposições~\ref{prop:envelope}--\ref{prop:bkidentity}, provadas
abaixo.
\end{corollary}

A construção de $H$, sua não-constância, e sua oscilação certificada estão na
Seção~\ref{sec:periodic}; a transferência de ponto de sela para $\varphi$ em si e a identificação
com a assintótica de Berg e Kr\"uppel, nas Seções~\ref{sec:formulaA}--\ref{sec:bk}.

\section{A densidade invariante e sua transformada de Laplace}
\label{sec:prelim}

Sejam $U_1, U_2, \dots$ variáveis aleatórias uniformes independentes e identicamente distribuídas
em $[0,1]$, e defina
\[
X = \sum_{j \ge 1} \frac{2}{3^j} U_j \in [0,1].
\]
A identidade autossimilar $X = \tfrac{2}{3}U_1 + \tfrac{1}{3}X'$, com $X'$ uma cópia independente de
$X$, mostra que a densidade $\varphi$ de $X$ satisfaz
\[
\varphi(x) = \frac{3}{2}\int_{3x-2}^{3x} \varphi,
\]
equivalentemente $\varphi'(x) = \tfrac{9}{2}\varphi(3x)$ em $(0, \tfrac{2}{3})$: a equação integral
exibida é exatamente a equação~(6.1) de \cite{wirsching2003}, $(W_3f)(x) := \tfrac32\int_{3x-2}^{3x}
f$, e $\varphi$ é exatamente a função que seu Corolário~7 constrói como a única solução
$L^1_{\rm loc}$ de $\operatorname{supp}\varphi\subset[0,1]$, $\int_0^1\varphi=1$, $W_3\varphi=\varphi$,
sua densidade invariante, a menos da reflexão $x \mapsto 1-x$ que fixa a convenção. É também uma
translação da função atômica $h_3$ de Rvachev \cite{rvachev1990}, e, em parâmetro de dilatação
$a=3$ e $\lambda=2/3$, uma instância da família de equações funcionais de Berg e Kr\"uppel
$\lambda\, \psi'(t) = a\bigl(\psi(at) - \psi(at-2)\bigr)$. A reflexão é de fato vazia,
$\varphi(x)=\varphi(1-x)$ identicamente: $X$ e $1-X=\sum_{j\ge1}\tfrac{2}{3^j}(1-U_j)$ são
equidistribuídas, já que $1-U_j\sim U_j$ independentemente para todo $j$ e
$\sum_{j\ge1}2\cdot3^{-j}=1$. Ela é suportada em $[0,1]$, contínua (pelo Lema~\ref{lem:tail}
abaixo), e constante, igual a $3/2$, no terço médio $[1/3, 2/3]$: ali, $3x-2\le0\le1\le3x$, de modo
que $\operatorname{supp}\varphi\subset[0,1]$ colapsa a integral da equação funcional a
$\int_0^1\varphi=1$, dando $\varphi(x)=\tfrac32\cdot1=\tfrac32$.

Para $s > 0$ escreva
\[
K(s) = \log \mathbb{E}\bigl[e^{-sX}\bigr] = \sum_{j \ge 1} g\!\left(\frac{2s}{3^j}\right), \qquad
g(b) = \log\frac{1-e^{-b}}{b}.
\]
O ramo de $g$ é o que é analítico em $\{\operatorname{Re} w > 0\}$ e real em $(0,\infty)$; já que
$(1-e^{-w})/w$ é analítica e não se anula ali (seus zeros estão em $2\pi i \Z \setminus \{0\}$, e
$w=0$ é uma singularidade removível de valor $1$), e o semiplano é simplesmente conexo, $g$ existe e
é única. Derivando termo a termo, justificado porque $g(w) = -w/2 + O(w^2)$ quando $w \to 0$ de modo
que a cauda da série é normalmente convergente junto com toda derivada,
\[
t(s) := -K'(s), \qquad V(s) := s^2 K''(s) > 0,
\]
a segunda uma positividade exata, sendo $K''(s)$ a variância de $X$ sob sua inclinação exponencial
por $e^{-sX}$. Como $t'(s) = -K''(s) < 0$, $t$ é estritamente decrescente em $(0,\infty)$; nos
extremos, $t(0^+) = \mathbb E[X] = \sum_{j\ge1} 2\cdot3^{-j}\cdot\tfrac12 = \tfrac12$, e $t(s)\to0$
quando $s\to\infty$ porque a inclinação exponencial se concentra em $\operatorname{ess\,inf} X=0$
(cada fator $e^{-sX}$ suprime qualquer massa longe de $0$ quando $s\to\infty$, e $X\ge0$ q.c.).
Assim $t(s)$ decresce estritamente de $1/2$ a $0$, e estudar $\varphi(t(s))$ quando $s \to \infty$ é
o mesmo que estudar $\varphi(t)$ quando $t \to 0^+$.

A função de comparação de Berg e Kr\"uppel é a solução assintótica da equação truncada
$\lambda\, \psi'(t) = a\, \psi(at)$, sendo $a$ aqui o parâmetro de dilatação deles (esta letra é
reutilizada na Seção~\ref{sec:periodic} para o coeficiente linear de $Q$, próprio deste artigo e
sem relação com aquele; os dois nunca aparecem na mesma fórmula exceto na Seção~\ref{sec:bk}, onde
a colisão é sinalizada de novo no ponto em que importa), por transformada de Laplace e análise de
ponto de sela. Na normalização deles (sua equação~(9.6), com
$\alpha = \tfrac12 - \tfrac{\log 2}{\log 3}$, $\beta = \tfrac{1}{2\log 3}$, e $\gamma$,
$\delta_{\mathrm{BK}}$, $\varepsilon$ determinados por $\alpha$ e $\beta$ através de suas próprias
equações, dadas explicitamente na Seção~\ref{sec:bk}),
\[
\varphi_0(t) = \frac{(2\beta)^\varepsilon}{\sqrt{2\pi}}\; t^\gamma\,(-\log t)^{\delta_{\mathrm{BK}}}
\, \exp\!\left(-\beta \log^2\!\frac{t}{-\log t}\right).
\]
Escreva $\varphi_{0,\mathrm{bare}}$ para a mesma expressão sem a constante inicial
$(2\beta)^\varepsilon/\sqrt{2\pi}$: as duas diferem apenas por essa constante.

\section{A correção periódica}
\label{sec:periodic}

Defina $c = \log 3$ e escreva $w = \log s$, $L(w) = K(e^w)$. Substituindo na série que define $K$ e
separando o termo $j=1$ do resto reproduz a recorrência
\begin{equation}
\label{eq:recurrence}
L(w+c) - L(w) = \log\bigl(1 - e^{-2e^w}\bigr) - w - \log 2.
\end{equation}
Seja
\[
Q(w) = -\frac{w^2}{2c} + \left(\frac12 - \frac{\log 2}{c}\right) w,
\]
escolhida de modo que $Q(w+c) - Q(w) = -w-\log2$, coincidindo exatamente com os termos
não-periódicos de \eqref{eq:recurrence}. Escreva $a := \tfrac12 - \tfrac{\log 2}{c}$ para o
coeficiente linear de $Q$, usado de novo a partir da Seção~\ref{sec:envelope}.

\begin{theorem}
\label{thm:H}
Defina
\begin{equation}
\label{eq:Hexact}
H(w) := L(w) - Q(w) + \sum_{k \ge 0} \log\bigl(1 - e^{-2\cdot 3^k e^w}\bigr).
\end{equation}
Então $H$ é exatamente $c$-periódica, $H(w+c) = H(w)$ para todo $w \in \R$, e
\[
L(w) = Q(w) + H(w) + \Delta(w), \qquad \Delta(w) = -\sum_{k\ge 0}\log\bigl(1-e^{-2\cdot3^k e^w}\bigr),
\]
com $\Delta^{(j)}(w) = O_j\bigl(e^{jw - 2e^w}\bigr)$ para todo $j \ge 0$: duplamente
exponencialmente pequeno, junto com todas as suas derivadas, quando $w \to +\infty$.
\end{theorem}

\begin{proof}
Escreva $d(w) = \log(1-e^{-2e^w})$, de modo que \eqref{eq:recurrence} se lê $L(w+c)-L(w) = d(w) - w
- \log 2$. Diretamente da definição de $Q$ (expanda $(w+c)^2 = w^2+2wc+c^2$ e agrupe os termos),
$Q(w+c)-Q(w) = -w-\log2$ também, de modo que as duas partes não-periódicas se cancelam exatamente:
\[
(L-Q)(w+c) - (L-Q)(w) = \bigl[d(w)-w-\log2\bigr] - \bigl[-w-\log2\bigr] = d(w).
\]
A série em \eqref{eq:Hexact} telescopia isso exatamente: com $D(w) := (L-Q)(w)$,
\[
D(w) + \sum_{k\ge0} d(w+kc) = D(w+c) - d(w) + \sum_{k\ge0}d(w+kc) = D(w+c) + \sum_{k\ge0}d(w+(k+1)c),
\]
de modo que o lado direito de \eqref{eq:Hexact} é invariante sob $w \mapsto w+c$, isto é, $H$ é
$c$-periódica. A convergência da série e sua diferenciabilidade termo a termo, junto com a cota
enunciada em $\Delta^{(j)}$, seguem de $d(w) = O(e^{-2e^w})$ e das cotas correspondentes de suas
derivadas, já que $2\cdot3^k e^w$ cresce como $3^k$ para $k\ge0$: os termos $d(w+kc)$ decaem
duplamente exponencialmente em $k$, portanto mais rápido do que qualquer taxa geométrica fixa, o
que dá convergência uniforme e dominada da série e de suas derivadas para $w$ em qualquer semirreta
limitada longe de $-\infty$.
\end{proof}

Dado $Q$, $H$ é unicamente determinada pela exigência de que $L-Q-H$ seja $o(1)$ quando $w\to\infty$:
duas funções periódicas quaisquer que diferem de $L-Q$ por um termo $o(1)$ quando $w\to\infty$ devem
ser iguais, sendo funções periódicas que coincidem no limite ao longo de toda classe residual de $w$
módulo $c$.

\begin{remark}
\label{rem:bkproduct}
A própria Proposição~9.3 de Berg e Kr\"uppel, especializada ao parâmetro de dilatação deles $3$ e
ao seu $b=3/2$ (exatamente esta autofunção; o $\alpha$ deles, mantido como símbolo aqui, é
exatamente o $\alpha$ importado da Seção~\ref{sec:prelim}, que a Proposição~\ref{prop:bkidentity}
abaixo identifica com o próprio $a$ deste artigo), dá a identidade equivalente
$\exp(H(w)) = 3^{\frac12 t^2 - \alpha t} \prod_{k\ge0}\frac{1-e^{-3^{t-k}/b}}{3^{t-k}/b}
\prod_{l\ge1}\bigl(1-e^{-3^{t+l}/b}\bigr)$, $t=w/c$, como um produto infinito, na prova daquela
proposição. Isto é $e^{-Q}e^Le^{-\Delta}$ termo a termo: escrevendo $s=e^w=3^t$, $3^{t-k}/b =
2s/3^{k+1}$, de modo que reindexar $j:=k+1\ge1$ transforma o primeiro produto em
$\prod_{j\ge1}\frac{1-e^{-2s/3^j}}{2s/3^j}$, exatamente $e^{L(w)}$ da série que define $K$;
$3^{t+l}/b = 2s\cdot3^{l-1}$, de modo que reindexar $k:=l-1\ge0$ transforma o segundo em
$\prod_{k\ge0}\bigl(1-e^{-2s\cdot3^k}\bigr)$, exatamente $e^{-\Delta(w)}$ por \eqref{eq:Hexact}; e,
já que $w=ct$ e $e^c=3$,
\[
-Q(w) = \frac{w^2}{2c}-aw = c\bigl(\tfrac12t^2-at\bigr), \qquad\text{de modo que}\qquad
3^{\frac12t^2-\alpha t} = e^{-Q(w)}
\]
uma vez que $\alpha=a$ (a Proposição~\ref{prop:bkidentity} identifica o $\alpha$ de Berg e
Kr\"uppel com o coeficiente linear de $Q$ deste artigo, também chamado $a$). Assim, o cálculo
primário da Proposição~\ref{prop:Hcert} a partir de \eqref{eq:Hcomputable} abaixo é, em essência,
uma avaliação desse produto. O que a série de Fourier da Proposição~\ref{prop:Hfourier} fornece, e
a forma de produto deles não, são as cotas de derivada \eqref{eq:Hbounds}, das quais dependem as
demais estimativas de ponto de sela do artigo.
\end{remark}

\begin{proposition}
\label{prop:Hfourier}
$H$ tem a expansão de Fourier $H(w) = \sum_{m\in\Z} \hat H(m) e^{i\omega_m w}$, $\omega_m =
2\pi m/c$, com
\[
\hat H(0) = \frac{\log 2}{2} - \frac{\log^2 2}{2c} - \frac{c}{12} - \frac{A}{c}, \qquad
A = \frac{\pi^2}{12} - \frac{\gamma_E^2}{2} - \gamma_1,
\]
\[
\hat H(m) = -\frac{2^{i\omega_m}}{c}\, \Gamma(-i\omega_m)\, \zeta(1-i\omega_m), \qquad m \ne 0,
\]
onde $\gamma_E$ é a constante de Euler--Mascheroni e $\gamma_1$ a primeira constante de Stieltjes.
\end{proposition}

\begin{proof}
Tanto $g$ quanto $K$ têm transformadas de Mellin na faixa $-1<\operatorname{Re}z<0$. Quando
$b\to0^+$, $g(b)=-b/2+O(b^2)$; quando $b\to\infty$, $g(b)=-\log b+O(e^{-b})$. Logo $\int_0^\infty
|g(b)|\,b^{\sigma-1}\,db<\infty$ para $-1<\sigma<0$, e as mesmas duas cotas aplicadas termo a termo
dão $K(s)=O(s)$ quando $s\to0^+$ e $K(s)=O(\log^2 s)$ quando $s\to\infty$: os termos com $3^j\le2s$
são em número $O(\log s)$ e cada um é $O(\log s)$, e os termos com $3^j>2s$ somam $O(1)$.

Integre por partes. Com $g'(b)=1/(e^b-1)-1/b$, e ambos os termos de contorno se anulando na faixa
($g(b)b^z\to0$ em $0$ porque $\operatorname{Re}z>-1$, e em $\infty$ porque $\operatorname{Re}z<0$),
\[
g^*(z) = \int_0^\infty g(b)\,b^{z-1}\,db = -\frac1z\int_0^\infty\Bigl(\frac{1}{e^b-1}-\frac1b\Bigr)b^z\,db
= -\frac{\Gamma(z+1)\zeta(z+1)}{z} = -\Gamma(z)\,\zeta(1+z),
\]
sendo o passo intermediário a continuação clássica
$\int_0^\infty\bigl(\tfrac{1}{e^u-1}-\tfrac1u\bigr)u^{v-1}\,du=\Gamma(v)\zeta(v)$, válida para
$0<\operatorname{Re}v<1$, em $v=z+1$.

Agora a soma harmônica. No $j$-ésimo termo substitua $u=2s/3^j$, de modo que $s=3^ju/2$ e
$s^{z-1}\,ds=(3^j/2)^z u^{z-1}\,du$:
\[
\int_0^\infty g\Bigl(\frac{2s}{3^j}\Bigr)s^{z-1}\,ds = \Bigl(\frac{3^j}{2}\Bigr)^{\!z}\int_0^\infty
g(u)\,u^{z-1}\,du = \Bigl(\frac{3^j}{2}\Bigr)^{\!z} g^*(z).
\]
Somar sobre $j\ge1$ é legítimo termo a termo, já que $\sum_{j\ge1}(3^j/2)^\sigma\int_0^\infty
|g(u)|u^{\sigma-1}\,du$ converge para $-1<\sigma<0$. O multiplicador que aparece é, portanto,
\[
\sum_{j\ge1}\Bigl(\frac{3^j}{2}\Bigr)^{\!z} = 2^{-z}\sum_{j\ge1}3^{jz} = \frac{2^{-z}3^z}{1-3^z}
= \frac{(3/2)^z}{1-3^z},
\]
uma série geométrica de razão $3^z$, de módulo $3^{\operatorname{Re}z}<1$, convergente em
$\operatorname{Re}z<0$ e em nenhum outro lugar. Logo
\begin{equation}
\label{eq:Kstar}
K^*(z) = \frac{(3/2)^z}{1-3^z}\cdot\bigl(-\Gamma(z)\zeta(1+z)\bigr)
= \frac{(3/2)^z\,\Gamma(z)\,\zeta(1+z)}{3^z-1}, \qquad -1<\operatorname{Re}z<0.
\end{equation}
As dilatações $\mu_j=2/3^j$ entram invertidas, através de $\mu_j^{-z}$, de onde vem o $2^{-z}$, e
com ele o $2^{i\omega_m}$ do enunciado.

Em $\operatorname{Re}z>-1$ o lado direito de \eqref{eq:Kstar} é meromorfo, com um polo de ordem
três em $z=0$ (cada um de $\Gamma(z)$, $\zeta(1+z)$ e $(3^z-1)^{-1}$ contribuindo um) e um polo
simples em cada zero restante $z=i\omega_m$, $m\ne0$, de $3^z-1$. À direita do eixo imaginário ela
é analítica.

Fixe $\sigma\in(-1,0)$ e $\sigma'\in(0,1)$, e desloque o contorno de inversão em
$K(s)=(2\pi i)^{-1}\int_{(\sigma)}K^*(z)s^{-z}\,dz$ de $\sigma$ para $\sigma'$, usando retângulos com
lados horizontais em $\operatorname{Im}z=\pm T_M$, $T_M:=\pi(2M+1)/c$. Nesses lados
$3^z=3^{\operatorname{Re}z}e^{i\pi(2M+1)}=-3^{\operatorname{Re}z}$, de modo que $|3^z-1| =
1+3^{\operatorname{Re}z}\ge1$; Stirling faz $|\Gamma(z)|$ decair como $e^{-\pi T_M/2}$ vezes uma
potência de $T_M$, uniformemente para $\sigma\le\operatorname{Re}z\le\sigma'$, e $\zeta(1+z)$ cresce
no máximo polinomialmente ali. As contribuições horizontais, portanto, se anulam quando
$M\to\infty$, a soma dos resíduos converge absolutamente, e na linha vertical direita
$|3^z-1|\ge3^{\sigma'}-1>0$ faz $\int_{(\sigma')}|K^*(z)s^{-z}|\,|dz| = O(s^{-\sigma'})$. Assim
\[
K(s) = -\operatorname*{Res}_{z=0}\bigl[K^*(z)s^{-z}\bigr]
-\sum_{m\ne0}\operatorname*{Res}_{z=i\omega_m}\bigl[K^*(z)s^{-z}\bigr]+O(s^{-\sigma'}).
\]

Tome primeiro os polos simples. Em $z=i\omega_m$ tem-se $3^{i\omega_m}=e^{2\pi im}=1$, de modo que
$(3^z-1)'=3^z\log3$ é igual a $c$ ali e $(3/2)^{i\omega_m}=3^{i\omega_m}2^{-i\omega_m}=
2^{-i\omega_m}$. Com $s=e^w$,
\[
-\operatorname*{Res}_{z=i\omega_m}\bigl[K^*(z)e^{-zw}\bigr]
= -\frac{2^{-i\omega_m}}{c}\,\Gamma(i\omega_m)\,\zeta(1+i\omega_m)\,e^{-i\omega_m w},
\]
e substituir $m$ por $-m$ em toda a soma transforma a família inteira em $\hat H(m)e^{i\omega_m w}$
com $\hat H(m)$ exatamente como enunciado.

Na origem, insira
\[
\Gamma(z) = \frac1z-\gamma_E+\Bigl(\frac{\gamma_E^2}{2}+\frac{\pi^2}{12}\Bigr)z+O(z^2), \qquad
\zeta(1+z) = \frac1z+\gamma_E-\gamma_1 z+O(z^2).
\]
Os dois termos $1/z$ se cancelam entre si e
\[
\Gamma(z)\zeta(1+z) = \frac{1}{z^2}+A+O(z), \qquad A = \frac{\pi^2}{12}-\frac{\gamma_E^2}{2}-\gamma_1,
\]
de onde vem $A$. Também $(3^z-1)^{-1} = (cz)^{-1}-\tfrac12+cz/12+O(z^3)$, e
$(3/2)^ze^{-zw}=e^{\theta z}$ com $\theta:=\log(3/2)-w$. Multiplicando as três expansões e lendo o
coeficiente de $z^{-1}$,
\[
\operatorname*{Res}_{z=0}\bigl[K^*(z)e^{-zw}\bigr]
= \frac1c\cdot\frac{\theta^2}{2}-\frac{\theta}{2}+\frac{c}{12}+\frac{A}{c}.
\]
Substitua $\theta=\log(3/2)-w$ e $\log(3/2)=c-\log2$. A parte de $-\operatorname{Res}$ que depende
de $w$ é
\[
-\frac{w^2}{2c}+\Bigl(\frac{\log(3/2)}{c}-\frac12\Bigr)w
= -\frac{w^2}{2c}+\Bigl(\frac12-\frac{\log2}{c}\Bigr)w = Q(w),
\]
e a parte constante é
$\tfrac12\log2-\tfrac{\log^22}{2c}-\tfrac{c}{12}-\tfrac Ac = \hat H(0)$.

Reunindo, com $s=e^w$,
\[
L(w) = Q(w)+\widetilde H(w)+O(e^{-\sigma' w}), \qquad
\widetilde H(w) := \hat H(0)+\sum_{m\ne0}\hat H(m)e^{i\omega_m w}.
\]
Essa série converge absoluta e uniformemente, porque $|\Gamma(-i\omega_m)| =
\sqrt{\pi/(\omega_m\sinh\pi\omega_m)}$ decai como $e^{-\pi\omega_m/2}$ enquanto $\zeta(1-i\omega_m)$
cresce polinomialmente, de modo que $\widetilde H$ é contínua e exatamente $c$-periódica. Ela é
real: $\Gamma$ e $\zeta$ são reais no eixo real, de modo que $\hat H(-m)$ é o conjugado complexo de
$\hat H(m)$.

O resto $O(e^{-\sigma'w})$ é apenas exponencialmente pequeno de forma ordinária, mais fraco do que
o duplamente exponencial $\Delta$ do Teorema~\ref{thm:H}, e isso já basta. O Teorema~\ref{thm:H} dá
$L-Q-H=\Delta=o(1)$ quando $w\to+\infty$, de modo que $\widetilde H-H = O(e^{-\sigma'w})+\Delta(w)\to0$.
Ambas são $c$-periódicas, e uma função $c$-periódica que tende a $0$ quando $w\to+\infty$ se anula
identicamente, sendo seu valor em qualquer $w$ igual ao seu valor em $w+kc$ para todo $k$. Logo
$\widetilde H=H$, e por convergência uniforme os $\hat H(m)$ são os coeficientes de Fourier de $H$.
\end{proof}

\begin{remark}
Isto é uma instância do fenômeno de de Bruijn--Mahler: uma soma harmônica de mudança de base,
transformada por Mellin, adquire polos simples numa linha vertical, e seus resíduos se reúnem numa
correção exatamente periódica com coeficientes de Fourier em $\Gamma$ e $\zeta$. De Bruijn encontrou
o primeiro caso disso na assintótica de partição de Mahler \cite{debruijn1948}; Erd\H{o}s e
Richmond o nomeiam na p.~448 de \cite{erdosrichmond1976}; Kato e McLeod \cite{katomcleod1971} e
Derfel \cite{derfel1995} estabeleceram a taxa de decaimento em log ao quadrado da assintótica de
comparação $\varphi_0$ da Seção~\ref{sec:prelim} como genérica em toda esta classe de equações
funcionais reescaladas. Nossa contribuição aqui é a forma fechada da própria correção periódica, não
o mecanismo que a produz.
\end{remark}

As séries de $H'$ e $H''$ convergem absoluta e uniformemente, já que
\[
|\Gamma(-i\omega_m)| = \sqrt{\pi/(\omega_m \sinh \pi\omega_m)}
\]
decai exponencialmente em $m$ enquanto $\zeta(1-i\omega_m)$ cresce apenas polinomialmente: para
$\omega\ge \alpha_H := 2\pi/c$ e $N:=\lfloor\omega\rfloor\ge1$, a fórmula clássica de
Euler--Maclaurin para $\zeta$, $\zeta(s) = \sum_{n=1}^N n^{-s} + \frac{N^{1-s}}{s-1} -
s\int_N^\infty\{u\}\,u^{-s-1}\,du$ (válida para $\operatorname{Re}s>0$, $s\ne1$, inteiro $N\ge1$,
com $\{u\}:=u-\lfloor u\rfloor$), em $s=1+i\omega$ dá, limitando cada termo em módulo
($|N^{-i\omega}|=1$, $0\le\{u\}<1$),
\[
|\zeta(1+i\omega)| \;\le\; \sum_{n=1}^N\frac1n + \frac1\omega + \frac{\sqrt{1+\omega^2}}{N}
\;\le\; 1+\log N + \frac1{\alpha_H} + \frac{\sqrt{1+\omega^2}}N,
\]
usando $\sum_{n=1}^N 1/n\le1+\log N$ e $\omega\ge\alpha_H$ no último termo. Como $N\le\omega$,
$\log N\le\log\omega$; como $N>\omega-1$ e $\omega\mapsto\omega/(\omega-1)$ é decrescente para
$\omega>1$, $\sqrt{1+\omega^2}/N < \bigl(\omega/(\omega-1)\bigr)\sqrt{1+\omega^{-2}} \le
\bigl(\alpha_H/(\alpha_H-1)\bigr)\sqrt{1+\alpha_H^{-2}}$. Com
\[
C_H := 1+\alpha_H^{-1}+\frac{\alpha_H}{\alpha_H-1}\sqrt{1+\alpha_H^{-2}},
\]
isso dá $|\zeta(1+i\omega)|\le\log\omega+C_H\le\log(\omega+1)+C_H$ para $\omega\ge\alpha_H$, e
$\zeta(1-i\omega)=\overline{\zeta(1+i\omega)}$ tem o mesmo módulo, de modo que
$|\zeta(1-i\omega_m)|\le\log(\omega_m+1)+C_H$. Com $q := e^{-\pi\alpha_H/2}$, $A_H :=
c^{-1}\sqrt{3\pi/\alpha_H}$, a cota $|\Gamma(-i\omega_m)|\le\sqrt{3\pi/\omega_m}\,e^{-\pi\omega_m/2}$
(válida para $\omega_m\ge\alpha_H$, já que $\sinh x\ge e^x/3$ ali) junto com a cota de
$\zeta(1-i\omega_m)$ recém-provada dão o majorante
\begin{equation}
\label{eq:Hmajorant}
|\hat H(m)| \le A_H\,q^m\,m^{-1/2}\bigl(\log(\alpha_H m+1)+C_H\bigr), \qquad m\ge1.
\end{equation}
Como $\sup_w|H^{(k)}(w)| \le \sum_{m\ne0}|\omega_m|^k|\hat H(m)| = 2\sum_{m\ge1}\omega_m^k|\hat
H(m)|$, somar os primeiros seis termos com envolvimentos rigorosos em bolas e limitar a cauda $m>6$
por \eqref{eq:Hmajorant} (a mesma soma de cauda geométrica-vezes-polinomial em forma fechada usada
na prova da Proposição~\ref{prop:Hcert} abaixo) dá cotas explícitas e certificadas:
\begin{equation}
\label{eq:Hbounds}
\sup_w |H'(w)| \le 0.0011977472315550332, \qquad \sup_w |H''(w)| \le 0.0068518962896650951.
\end{equation}

\begin{proposition}[Não-constância e oscilação certificada]
\label{prop:Hcert}
$H$ não é constante: $\hat H(1)\ne0$ segue do teorema clássico de não-anulação de $\zeta$ em
$\operatorname{Re}s=1$ sozinho, sem nenhum cálculo. A oscilação é quantificada rigorosamente por um
certificado explícito de aritmética de intervalos:
\[
-0.000377190280943987 \;<\; H(0) - H(\log(3/2)) \;<\; -0.000377190280943985,
\]
e consequentemente $\operatorname{osc}(H) := \sup H - \inf H \ge 3.7719\times 10^{-4}$; avaliar $H$
em dois pontos adicionais, bem separados, certifica a cota mais forte
$\operatorname{osc}(H) \ge 4.1874494771\times10^{-4}$, e o modo de Fourier líder sozinho certifica
uma cota superior correspondente, dando juntas
\[
4.1874494771\times10^{-4} \;\le\; \operatorname{osc}(H) \;\le\; 4.187981\times10^{-4}.
\]
\end{proposition}

\begin{proof}
Primeiro a não-constância, incondicionalmente. Pela Proposição~\ref{prop:Hfourier},
$\hat H(1) = -\frac{2^{i\omega_1}}{c}\Gamma(-i\omega_1)\zeta(1-i\omega_1)$, $\omega_1 = 2\pi/c$, um
produto de três fatores, cada um não-nulo: $2^{i\omega_1}\ne0$ trivialmente; $\Gamma$ não tem zeros
em lugar nenhum, e $-i\omega_1$ não é um de seus polos ($\omega_1\ne0$ é real, e os polos de
$\Gamma$ estão apenas nos inteiros não-positivos), de modo que $\Gamma(-i\omega_1)\ne0$; e
$1-i\omega_1$ está na reta $\operatorname{Re}s=1$ com $\omega_1\ne0$, onde $\zeta(1+it)\ne0$ para
todo real $t\ne0$ pelo teorema clássico de não-anulação de Hadamard e de la Vall\'ee Poussin
\cite{titchmarsh1986}, de modo que $\zeta(1-i\omega_1)\ne0$. Logo $\hat H(1)\ne0$, e uma função
periódica com um coeficiente de Fourier não-nulo em $m\ne0$ não pode ser constante.

Para a cota quantitativa, escreva $g(b) = -b/2 + \log S(b/2)$, $S(x) = \sinh(x)/x = \sum_{n\ge0}
x^{2n}/(2n+1)!$ (esta é a fórmula que define $g$, reescrita em torno de sua singularidade removível
em $0$, já que $e^{-b/2}(e^{b/2} - e^{-b/2}) = 1-e^{-b}$). A dominação termo a termo $6^n n! \le
(2n+1)!$, que vale em $n=0$ e se propaga porque $(2n+3)(2n+2)\ge6(n+1)$, dá $1\le S(x)\le e^{x^2/6}$,
isto é
\begin{equation}
\label{eq:logSbound}
0 \;\le\; \log S(x) \;\le\; \frac{x^2}{6}, \qquad x\ge0.
\end{equation}
Como $\sum_{j\ge1}3^{-j} = \tfrac12$, as metades $-b/2$ dos termos de $L$ somam em forma fechada, e
\eqref{eq:Hexact} se torna, com $s = e^w$,
\begin{equation}
\label{eq:Hcomputable}
H(w) = -\frac s2 + \sum_{j\ge1}\log S\!\left(\frac{s}{3^j}\right) - Q(w)
+ \sum_{k\ge0}\log\bigl(1-e^{-2\cdot3^k s}\bigr).
\end{equation}
Trunque a primeira série em $j=R$ e a segunda em $k=M$. Por \eqref{eq:logSbound} a primeira cauda
satisfaz
\[
0 \;\le\; \sum_{j>R}\log S\!\left(\frac{s}{3^j}\right) \;\le\; \frac{s^2}{6}\sum_{j>R}9^{-j}
= \frac{s^2\,9^{-R}}{48},
\]
e, como $0\le-\log(1-y)\le y/(1-y)$ e $e^{-2\cdot3^ks} =
\bigl(e^{-2\cdot3^{M+1}s}\bigr)^{3^{k-M-1}}$ para $k>M$, a segunda cauda satisfaz
\[
0 \;\le\; -\sum_{k>M}\log\bigl(1-e^{-2\cdot3^ks}\bigr) \;\le\; \frac{y}{(1-y)^2} \;\le\; 4y,
\qquad y:=e^{-2\cdot3^{M+1}s}\le\tfrac12 .
\]
Tome $R=30$, $M=4$. Para $s\le3$ (o que cobre todo ponto de avaliação usado abaixo, incluindo a
grade de oscilação) a primeira cauda é no máximo $4.43\times10^{-30}$, e para $s\ge1$ a segunda é
no máximo $4e^{-486}<10^{-210}$. Os $35$ termos restantes são funções elementares de $s$; avaliados
com $50$ dígitos decimais significativos, onde o arredondamento acumulado permanece abaixo de
$10^{-45}$, eles dão
\[
H(0) = -0.62713093305152597830\ldots, \qquad H\bigl(\log\tfrac32\bigr) = -0.62675374277058199247\ldots,
\]
cuja diferença é $-0.00037719028094398583\ldots$, dentro do intervalo enunciado; esta primeira rota
é em ponto flutuante ($50$ dígitos significativos, arredondamento acumulado limitado como
enunciado) em vez de aritmética de intervalos. A série de Fourier da Proposição~\ref{prop:Hfourier}
dá a mesma diferença em aritmética de bolas de Arb, rigorosamente: o termo $m=0$ se cancela, e em
precisão de trabalho de $250$ bits (raio de bola $<10^{-70}$ ao longo de todo o cálculo) os quatro
modos não-nulos necessários são
\begin{align*}
\hat H(1) &= -1.020544273068431785\times10^{-4} - 2.332597606471452884\times10^{-5}\,i,\\
\hat H(2) &= -6.523276079851823231\times10^{-9} + 1.156720838500605667\times10^{-8}\,i,\\
\hat H(3) &= -1.799881020474424677\times10^{-12} + 6.648927223521397737\times10^{-13}\,i,\\
\hat H(4) &= -9.219524659678388767\times10^{-17} + 9.444517590655925766\times10^{-17}\,i,
\end{align*}
a soma dos quatro modos
\[
2\operatorname{Re}\sum_{m=1}^4\hat H(m)\bigl(1-e^{i\omega_m\log(3/2)}\bigr)
= -0.0003771902809439858148\ldots
\]
(raio de bola $<10^{-70}$), coincidindo com o valor em ponto flutuante acima em $19$ casas decimais
($16$ dígitos significativos). O majorante \eqref{eq:Hmajorant} dá
$\sum_{m>4}|\hat H(m)| \le 9.41\times10^{-20}$, de modo que os modos descartados deslocam a soma dos
quatro modos por menos de $3.8\times10^{-19}$, dando
\[
H(0)-H\bigl(\log\tfrac32\bigr) \in \bigl[-0.0003771902809439861948,\; -0.0003771902809439854348\bigr];
\]
esta rota sozinha certifica o intervalo enunciado para $H(0)-H(\log(3/2))$.

Para a cota mais forte, seja $w_i = ic/N$, $0\le i<N$, com $N = 2^{20}$. Uma avaliação em ponto
flutuante dos $N$ valores da grade a partir da Proposição~\ref{prop:Hfourier}, truncada em
$|m|\le10$ (o mesmo majorante limita o erro de truncamento pontual em cada ponto por
$2\times2.80\times10^{-43}=5.60\times10^{-43}$, combinando os modos descartados $m$ e $-m$), localiza
um candidato a máximo em $i=486746$ e um candidato a mínimo em $i=1011118$; esta busca é apenas uma
heurística para encontrar dois pontos bem separados e não faz parte, ela mesma, do certificado
abaixo. Reavaliando $H$ exatamente nesses dois pontos com precisão de trabalho de $250$ bits em
aritmética de bolas de Arb, independentemente, a partir de \eqref{eq:Hcomputable} e da série de
Fourier (as duas rotas coincidem a cerca de $30$ dígitos, consistente com a própria cota de cauda de
truncamento $R=30$ de \eqref{eq:Hcomputable} acima; uma verificação cruzada, não parte do
certificado),
\[
H(w_{486746}) = -0.6267174403736425786\ldots, \qquad H(w_{1011118}) = -0.6271361853213578093\ldots,
\]
de modo que $D := H(w_{486746})-H(w_{1011118}) = 4.187449477152\times10^{-4}$. O envolvimento
certificado usa apenas a rota de Fourier: o arredondamento de bola de Arb em precisão de $250$
bits é $<10^{-70}$ por avaliação, desprezível diante do erro de truncamento do majorante já
limitado acima, $5.60\times10^{-43}$ por ponto, dando $D$ com raio $<1.2\times10^{-42}$ (o dobro
dessa cota por ponto). Já que $w_{486746}$ e $w_{1011118}$ são dois pontos específicos e fixos do
domínio, $H(w_{486746})\le\sup H$ e $H(w_{1011118})\ge\inf H$ valem trivialmente, dando
$D\le\operatorname{osc}(H)$: essa cota é incondicional, independente de a busca em ponto flutuante
acima ter de fato localizado o verdadeiro argmax e argmin, já que ela usa apenas que os dois pontos
escolhidos estão no domínio.

Para a cota superior correspondente, nenhuma busca é necessária. Escreva $H = \hat H(0) +
2\operatorname{Re}\bigl(\hat H(1)e^{i\omega_1 w}\bigr) + \Xi(w)$, $\Xi(w) :=
2\sum_{m\ge2}\operatorname{Re}\bigl(\hat H(m) e^{i\omega_m w}\bigr)$, de modo que
$|\Xi(w)|\le2\sum_{m\ge2}|\hat H(m)|$ para todo $w$. O termo de um modo
$2\operatorname{Re}(\hat H(1)e^{i\omega_1w})$ varia, quando $w$ varia, sobre um intervalo de
largura exatamente $4|\hat H(1)|$, de modo que $\operatorname{osc}(H) \le 4|\hat H(1)| +
4\sum_{m\ge2}|\hat H(m)|$. Usando os já certificados $\hat H(2)$, $\hat H(3)$, $\hat H(4)$ de acima
e a cauda do majorante $\sum_{m>4}|\hat H(m)|\le9.41\times10^{-20}$,
\[
\sum_{m\ge2}|\hat H(m)| \le |\hat H(2)|+|\hat H(3)|+|\hat H(4)|+9.41\times10^{-20} \le
1.3282\times10^{-8},
\]
e $4|\hat H(1)| \le 4.187449304\times10^{-4}$, dando $\operatorname{osc}(H) \le
4.187981\times10^{-4}$, rigorosa e incondicionalmente: esta cota usa apenas os coeficientes de
Fourier já certificados, não a busca em grade.
\end{proof}

\section{Uma assintótica uniforme de ponto de sela}
\label{sec:formulaA}

Defina, para $s>0$,
\[
\varphi_{\mathrm{sp}}(t(s)) := \frac{\exp\bigl(K(s)+st(s)\bigr)}{\sqrt{2\pi K''(s)}}.
\]
Esta seção prova que $\varphi(t(s))/\varphi_{\mathrm{sp}}(t(s)) \to 1$ quando $s \to \infty$,
uniformemente sobre $\rho:=2s/3^N\in[1,3)$ (definido abaixo; este parâmetro registra a mesma fase,
$w\bmod c$, rastreada através de $H$ nas Seções~\ref{sec:bk} e~\ref{sec:proofs}). É essa uniformidade
que permite que os dois teoremas da Seção~\ref{sec:proofs} dependam de escopos diferentes em vez de
matemática diferente.

Para $n \ge 2$, defina as funções normalizadas por derivada $\kappa_n(w) = (-1)^n w^n g^{(n)}(w)$ em
$\{\operatorname{Re} w > 0\}$; $\kappa_n(b)$, $b>0$, é $b^n$ vezes o $n$-ésimo cumulante da variável
aleatória com densidade $b\,e^{-bv}/(1-e^{-b})$ em $[0,1]$. Escrevendo $x=w/2$,
\[
\kappa_2(w) = 1 - \frac{x^2}{\sinh^2 x}, \qquad
\kappa_3(w) = 2 - \frac{2x^3\cosh x}{\sinh^3 x}, \qquad
\kappa_4(w) = 6 + \frac{2x^4(1-3\coth^2 x)}{\sinh^2 x},
\]
e a partir da expansão de Bernoulli $g(w)+w/2 = \sum_{k\ge1} B_{2k}w^{2k}/(2k(2k)!)$ em $|w|<2\pi$,
\[
\kappa_n(w) = (-1)^n\!\!\sum_{k\ge\lceil n/2\rceil}\!\! \frac{B_{2k}\,w^{2k}}{2k\,(2k-n)!}, \qquad n\ge2,
\]
de modo que $\kappa_n(w) = O(w^2)$ quando $w \to 0$ para todo $n \ge 2$.

Escreva $s = \rho\cdot 3^N/2$ para o único $N = \lfloor \log_3(2s)\rfloor \in \Z$ e $\rho \in
[1,3)$, e $b_j = 2s/3^j$. Então $\{b_j : j \ge 1\} = \{\rho\cdot 3^m : m \le N-1\}$; para $N \ge 1$
(o único intervalo usado a partir daqui, correspondendo a $s \ge 3/2$), exatamente $N$ desses, o
menor sendo $\rho \ge 1$, são $\ge 1$, e o resto forma uma cauda geométrica abaixo de $1$. Este único
fato estrutural é a origem de toda constante independente de fase abaixo.

\begin{lemma}[Limitação fora do eixo real]
\label{lem:kappa}
Para $w = b(1+iu)$, $b>0$, $|u|\le 1$, com $e(r) = r^2/\sinh^2 r$ e $f(r) = r^3\cosh r/\sinh^3 r$,
\[
|\kappa_2(w)| \le 1 + 2e(b/2), \qquad |\kappa_3(w)| \le 2 + 2\cdot 2^{3/2} f(b/2).
\]
Ambas as cotas são finitas em todo o setor $|\operatorname{Im}w|\le\operatorname{Re}w$:
$\sup|\kappa_2|\le3$, $\sup|\kappa_3|\le2+2^{5/2}<7.657$. Além disso, para $|w|\le2$,
$|\kappa_2(w)| \le 0.114|w|^2$ e $|\kappa_3(w)|\le 0.0119|w|^4$.
\end{lemma}

\begin{proof}
$|\sinh z|^2 = \sinh^2(\operatorname{Re}z)+\sin^2(\operatorname{Im}z) \ge \sinh^2(\operatorname{Re}z)$
e $|\cosh z|^2 \le \cosh^2(\operatorname{Re}z)$; a primeira dá $|\sinh(w/2)|\ge\sinh(b/2)>0$,
excluindo os polos de $\coth$, já que $\operatorname{Re}(w/2)=b/2>0$. Isso dá as cotas de setor
diretamente. A monotonicidade de $e$ é imediata da série de potências crescente de $\sinh(r)/r$.
Para $f$, ponha $u=2r$; um cálculo direto dá $2r\sinh(r)\cosh(r)\,(\log f)'(r) = F(u)$ com
$F(u):=3\sinh u - u\cosh u - 2u$, de modo que $(\log f)'(r)$ tem o sinal de $F(2r)$. Como
$F(0)=F'(0)=F''(0)=0$ e $F'''(u)=-u\sinh u<0$ para $u>0$, integrar três vezes a partir de $0$ dá
$F(u)<0$ para todo $u>0$, logo $(\log f)'(r)<0$ para $r>0$ e $f$ é estritamente decrescente.
(A desigualdade de Lazarevi\'c, $(\sinh r/r)^3 > \cosh r$, dá separadamente $f<1$.) As cotas locais
em $|w|\le2$ seguem aplicando o princípio do módulo máximo a $\kappa_2(w)/w^2$ e $\kappa_3(w)/w^4$
(analíticas em $|w|<2\pi$ pela expansão de Bernoulli): limitando cada uma em $|w|=2$ pela
desigualdade triangular, usando $|B_{2k}| = 2(2k)!\,\zeta(2k)/(2\pi)^{2k}$ e $\zeta(2k)\le\zeta(2)$,
respectivamente $\zeta(4)$.
\end{proof}

\begin{lemma}[Variância]
\label{lem:variance}
Para todo $\rho\in[1,3)$, $N\ge1$, com $V := s^2K''(s) = \sum_j\kappa_2(b_j)$,
\[
N - A_V \le V \le N + 0.1283, \qquad A_V := \sum_{m\ge0} e(3^m/2) = 1.4269413069\ldots.
\]
\end{lemma}

\begin{proof}
Separe $V = \sum_{m=0}^{N-1}\kappa_2(\rho3^m) + \sum_{m<0}\kappa_2(\rho3^m)$, usando $\kappa_2(b) =
1-e(b/2) \in (0,1)$. Para a cota inferior, descarte a segunda soma (não-negativa) e use
$e(\rho3^m/2) \le e(3^m/2)$ ($e$ decrescente, $\rho\ge1$) na primeira, dando
$\sum_{m=0}^{N-1}\kappa_2(\rho3^m) \ge N - A_V$. Para a cota superior, a primeira soma é $<N$ termo
a termo, e a segunda soma, onde $\rho3^m<1$, é limitada pela estimativa local do Lema~\ref{lem:kappa}:
$\sum_{m<0}\kappa_2(\rho3^m) \le 0.114\rho^2\sum_{m<0}9^m = 0.114\rho^2/8 \le 0.1283$ para $\rho<3$.
\end{proof}

\begin{lemma}[Cota de cauda e existência da densidade]
\label{lem:tail}
Se $b_0\ge1$ e $b_k=3^kb_0$, então
\[
\sum_{k\ge0}\log\coth(b_k/2) \le 3e^{-b_0}.
\]
Consequentemente, com $\Lambda(y) := K(s(1+iy)) - K(s) + isy\,t(s)$,
\[
|e^{\Lambda(y)}| \le e^{3/e}(1+y^2)^{-N/2} \le 3.0152\,(1+y^2)^{-N/2}
\]
para todo real $y$. Em particular $X$ tem densidade contínua, a função característica inclinada é
integrável para $N\ge3$, e a inversão de Fourier em $x=t(s)$ dá
\begin{equation}
\label{eq:R}
R(s) := \frac{\varphi(t(s))}{\varphi_{\mathrm{sp}}(t(s))} = \sqrt{\frac{V}{2\pi}}\int_\R e^{\Lambda(y)}\,dy,
\end{equation}
uma quantidade real, já que $\Lambda(-y) = \overline{\Lambda(y)}$.
\end{lemma}

\begin{proof}
$\log\coth(x/2) = 2\artanh(e^{-x}) \le 2e^{-x}/(1-e^{-2x})$. Escreva $q:=e^{-b_0}\in(0,e^{-1}]$, de
modo que $e^{-b_k}=q^{3^k}$ e a cota se lê $\log\coth(b_k/2) \le 2q^{3^k}/(1-q^{2\cdot3^k})$. Como
$3^k\ge1$ para $k\ge0$, $q^{2\cdot3^k}\le q^2$ termo a termo, de modo que o denominador é limitado
inferiormente por $1-q^2$ uniformemente em $k$:
\[
\sum_{k\ge0}\log\coth(b_k/2) \le \frac{2}{1-q^2}\sum_{k\ge0}q^{3^k}.
\]
Como $3^k\ge2k+1$ para todo $k\ge0$ e $0<q<1$, $q^{3^k}\le q^{2k+1}$, de modo que
$\sum_{k\ge0}q^{3^k} \le q\sum_{k\ge0}q^{2k} = q/(1-q^2)$, dando
$\sum_{k\ge0}\log\coth(b_k/2) \le 2q/(1-q^2)^2$. Para $b_0\ge1$, $q\le e^{-1}<0.4283$, de modo que
$(1-q^2)^2>2/3$ e $2q/(1-q^2)^2<3q=3e^{-b_0}$.
A densidade de $X$ é contínua porque a soma dos dois primeiros termos já tem densidade contínua e
de suporte compacto e a convolução preserva continuidade. Para $\Lambda$: $t(s) = -K'(s)$ cancela
exatamente o termo linear da transformada inclinada, e cada fator restante
$e^{g(b_j(1+iy))-g(b_j)}$ tem módulo no máximo $1$ (uma função característica inclinada), com módulo
no máximo $\coth(b_j/2)/\sqrt{1+y^2}$ nos $N$ índices onde $b_j\ge1$; a estimativa de cauda se aplica
exatamente a esse conjunto geométrico, com $b_0=\rho\ge1$.
\end{proof}

\begin{lemma}[Expansão central]
\label{lem:taylor}
Com $h(u) = g(b(1+iu))$, o teorema de Taylor na forma \emph{integral} (a forma de Lagrange é
inválida para uma função de valor complexo de uma variável real) dá
\[
h(y) = h(0)+h'(0)y+h''(0)\tfrac{y^2}{2} + \tfrac12\int_0^y (y-u)^2 h'''(u)\,du.
\]
Consequentemente, escrevendo $h_j(u) := g(b_j(1+iu))$ para cada $j\ge1$, para $\theta\in(0,1]$,
\[
S := \sum_j \sup_{|u|\le \theta}|h_j'''(u)| \le 2N + 10.559,
\]
e portanto $|\Lambda(y) + Vy^2/2| \le B|y|^3$ para $|y|\le \theta$, com $B := S/6 \le (2N+10.559)/6$.
\end{lemma}

\begin{proof}
Escreva $w_j := b_j(1+iu)$. A regra da cadeia e a definição de $\kappa_3$ dão $h_j'''(u) = i\,
\kappa_3(w_j)/(1+iu)^3$, de modo que
\[
|h_j'''(u)| = \frac{|\kappa_3(w_j)|}{(1+u^2)^{3/2}}.
\]
Separe o conjunto de índices como na Seção~\ref{sec:formulaA}. Para os $N$ índices com
$b_j = \rho3^m \ge 1$, $0\le m\le N-1$, descarte o denominador e use a cota de setor do
Lema~\ref{lem:kappa}, legítima já que $|u|\le\theta\le1$:
\[
|h_j'''(u)| \le |\kappa_3(w_j)| \le 2 + 2^{5/2} f(b_j/2) \le 2 + 2^{5/2} f(3^m/2),
\]
o último passo pela monotonicidade de $f$ e $\rho\ge1$. Somando sobre $m\ge0$ limita esse bloco por
$2N + 2^{5/2}\Sigma$, $\Sigma := \sum_{m\ge0} f(3^m/2)$. Os três primeiros termos de $\Sigma$ são
$0.99614738\ldots$, $0.82241066\ldots$, $0.04500508\ldots$; para $r\ge13$ tem-se $f(r) \le
4.001\,r^3e^{-2r}$, de modo que os termos restantes somam menos que $2\times10^{-8}$, e
$2^{5/2}\Sigma < 10.541906$.

Para os índices restantes $b_j = \rho 3^{-n}$, $n\ge1$, tem-se $b_j<1$ e $|w_j| =
b_j\sqrt{1+u^2} \le b_j\sqrt2 < 2$, de modo que a cota local do Lema~\ref{lem:kappa} se aplica e,
mantendo o denominador desta vez,
\[
|h_j'''(u)| \le \frac{0.0119\,|w_j|^4}{(1+u^2)^{3/2}} = 0.0119\,b_j^4\,(1+u^2)^{1/2} \le
0.0119\sqrt2\,b_j^4 .
\]
Somando, $\sum_{n\ge1}(\rho3^{-n})^4 = \rho^4/80 < 81/80$, de modo que esse bloco fica abaixo de
$0.0119\sqrt2\cdot81/80 < 0.01704$. Juntando, $S < 2N + 10.541906 + 0.01704 < 2N + 10.559$.

Para a cota do resto, $\sum_j h_j(y) = K(s(1+iy))$, $\sum_j h_j'(0) = isK'(s) = -is\,t(s)$ e
$\sum_j h_j''(0) = -s^2K''(s) = -V$, de modo que somar a expansão de Taylor em forma integral sobre
$j$ dá
\[
\Lambda(y) + \frac{Vy^2}{2} = \sum_j \frac12\int_0^y (y-u)^2 h_j'''(u)\,du,
\]
cujo módulo é no máximo $S|y|^3/6 = B|y|^3$ para $|y|\le\theta$.
\end{proof}

\begin{theorem}[Assintótica uniforme de ponto de sela]
\label{thm:formulaA}
Para $N\ge19$, uniformemente sobre $\rho\in[1,3)$, $|R(s)-1|\le E(N) := e_1(N)+e_2(N)+e_3(N)$, com
$\theta_N := N^{-1/3}$, $V_{\mathrm{lo}} := N-A_V$, $V_{\mathrm{up}} := N+0.1283$,
$B := (2N+10.559)/6$,
\begin{align}
e_1(N) &:= \frac{4\,e^{B\theta_N^3}\,B}{\sqrt{2\pi}\,V_{\mathrm{lo}}^{3/2}}, \qquad
e_2(N) := \frac{2}{\sqrt{2\pi}}\cdot\frac{e^{-\theta_N^2V_{\mathrm{lo}}/2}}{\theta_N\sqrt{V_{\mathrm{lo}}}},
\label{eq:Ebound}\\
e_3(N) &:= 2\,e^{3/e}\sqrt{\frac{V_{\mathrm{up}}}{2\pi}}\cdot
\frac{(1+\theta_N^2)^{-(N-2)/2}}{\theta_N(N-2)}. \notag
\end{align}
$\sqrt N\,E(N) \to \tfrac43e^{1/3}(2\pi)^{-1/2} = 0.742358\ldots$ quando $N\to\infty$
(provado abaixo, a partir de $e_1$ sozinho: $e_2$ e $e_3$ se anulam mais rápido do que qualquer
potência de $N$). Consequentemente $\varphi(t) = \varphi_{\mathrm{sp}}(t)\,(1+o(1))$ quando
$t \to 0^+$, sem nenhuma restrição sobre a sequência de $t$'s.
\end{theorem}

\begin{remark}
A estimativa $|R(s)-1|\le E(N)$ se torna informativa uma vez que $E(N)<1$, o que vale em $N=19$
(veja a prova: $E(18)=1.00170\ldots\ge1$, $E(19)=0.95673\ldots<1$); esse é o limiar mantido como
hipótese do teorema (se algum $N$ menor também dá $E(N)<1$ é irrelevante para o que segue e não é
verificado). $E$ também é estritamente decrescente e satisfaz $E(N)\le4.18\,N^{-1/2}$ ao longo de
$19\le N\le5000$ (verificado diretamente no repositório que acompanha o artigo), embora apenas
$E(N)\to0$ quando $N\to\infty$ seja usado abaixo.
\end{remark}

\begin{proof}
Separe \eqref{eq:R} em $|y|=\theta_N$. Em $|y|\le\theta_N$, o Lema~\ref{lem:taylor} limita o resto
cúbico $|\Lambda(y)+Vy^2/2|$ por $B\theta_N^3$, que permanece limitado (não $o(1)$: é exatamente por
isso que $e_1$ carrega o fator $e^{B\theta_N^3}$ em vez de se anular) quando $N\to\infty$, já que
$B=O(N)$ e $\theta_N^3=N^{-1}$. Escreva $C(y):=e^{\Lambda(y)}-e^{-Vy^2/2}$ e fatore a gaussiana antes
de estimar, $e^{\Lambda(y)} = e^{-Vy^2/2}\,e^{z(y)}$ com $z(y) := \Lambda(y)+Vy^2/2$, de modo que
$C(y) = e^{-Vy^2/2}\bigl(e^{z(y)}-1\bigr)$. Como $|z(y)|\le B|y|^3 \le B\theta_N^3$ em
$|y|\le\theta_N$, a desigualdade elementar $|e^z-1|\le|z|e^{|z|}$ dá
\[
|C(y)| \;\le\; e^{-Vy^2/2}\,B|y|^3\,e^{B\theta_N^3} \qquad (|y|\le\theta_N),
\]
mantendo o fator gaussiano. Integrando essa desigualdade, $\int_\R e^{-Vy^2/2}|y|^3\,dy=4/V^2$, de
modo que
\[
\sqrt{V/2\pi}\int_{|y|\le\theta_N}|C(y)|\,dy \;\le\; \frac{4Be^{B\theta_N^3}}{\sqrt{2\pi}\,V^{3/2}},
\]
que é decrescente em $V$; usar $V\ge V_{\mathrm{lo}}$ do Lema~\ref{lem:variance} dá a contribuição
central do resto de Taylor $e_1$. Como $\sqrt{V/2\pi}\int_\R e^{-Vy^2/2}\,dy = 1$, essa mesma divisão
deixa exatamente dois termos a mais. A cauda gaussiana truncada
$\sqrt{V/2\pi}\cdot2\int_{\theta_N}^\infty e^{-Vy^2/2}\,dy$, substituindo $v=y\sqrt V$ para obter
$(2/\sqrt{2\pi})\int_{c_0}^\infty e^{-v^2/2}\,dv$, $c_0=\theta_N\sqrt V$ (uma constante local,
distinta de $c:=\log3$), e limitando via $\int_{c_0}^\infty e^{-v^2/2}\,dv\le e^{-c_0^2/2}/c_0$ em
$V=V_{\mathrm{lo}}$, dá $e_2$. Em $|y|>\theta_N$, a cota do Lema~\ref{lem:tail}
$|e^{\Lambda(y)}|\le e^{3/e}(1+y^2)^{-N/2}$ se integra, via $\int_u^\infty(1+y^2)^{-N/2}\,dy \le
(1+u^2)^{-(N-2)/2}/(u(N-2))$ em $u=\theta_N$ (usando $V\le V_{\mathrm{up}}$ para o fator principal
$\sqrt{V/2\pi}$ em \eqref{eq:R}), para dar $e_3$; isso precisa de $\theta_N \gg N^{-1/2}$, satisfeito
já que $\tfrac13<\tfrac12$. Toda constante que entra em $E$ (via $A_V$, Lema~\ref{lem:variance}, e o
$3.0152\ge e^{3/e}$ do Lema~\ref{lem:tail}) é um supremo sobre $b>0$ ou uma soma sobre uma órbita
geométrica completa limitada usando $\rho\ge1$, portanto independente de $\rho$. Avaliar
diretamente \eqref{eq:Ebound} dá $E(18) = 1.00170\ldots \ge 1$ e $E(19) = 0.95673\ldots < 1$,
coincidindo com a observação acima.

Para a taxa assintótica, $e_1$, $e_2$, $e_3$ se separam em escalas muito diferentes quando
$N\to\infty$. Como $\theta_N^3=1/N$, $B\theta_N^3\to1/3$ e $B/N\to1/3$,
$V_{\mathrm{lo}}/N,\,V_{\mathrm{up}}/N \to1$, de modo que
\[
\sqrt N\,e_1(N) = \frac{4e^{B\theta_N^3}}{\sqrt{2\pi}}\cdot\frac{\sqrt N\,B}{V_{\mathrm{lo}}^{3/2}}
\longrightarrow \frac{4e^{1/3}}{\sqrt{2\pi}}\cdot\frac13 = \frac43e^{1/3}(2\pi)^{-1/2}.
\]
Para $e_2$ e $e_3$, $\theta_N^2V_{\mathrm{lo}} \sim N^{1/3}\to\infty$, de modo que ambos carregam um
fator $e^{-N^{1/3}/2}(1+o(1))$ (para $e_3$, via $(1+\theta_N^2)^{-(N-2)/2} =
\exp(-\tfrac{N-2}2\theta_N^2(1+o(1)))$), contra prefatores polinomiais de ordem $N^{-1/6}$; um fator
que decai como $e^{-N^{1/3}/2}$ vence toda potência de $N$, de modo que $N^k e_2(N),\,N^k
e_3(N)\to0$ para todo $k$. Logo $\sqrt N\,E(N) = \sqrt N\,e_1(N) + o(1) \to
\tfrac43e^{1/3}(2\pi)^{-1/2}$, e em particular $E(N)\to0$. Como $t(s)$ percorre bijetivamente
$(0,1/2)$, isso é exatamente o enunciado $\varphi(t)=\varphi_{\mathrm{sp}}(t)(1+o(1))$ quando
$t\to0^+$.
\end{proof}

\begin{remark}
\label{rem:edgeworth}
A taxa não é ótima, e uma mais fina é plausível mas não é estabelecida aqui. Levar a expansão
central à quarta ordem, usando que o termo cúbico ímpar se integra a zero contra o núcleo
gaussiano, e usando os limites $V_3 := \sum_j\kappa_3(b_j) = 2N+O(1)$, $V_4 :=
\sum_j\kappa_4(b_j) = 6N+O(1)$ (as somas do terceiro e quarto cumulantes, pelo mesmo argumento do
Lema~\ref{lem:variance}) sugere a correção de Edgeworth $R(s) = 1 - \tfrac1{12N} + O(N^{-3/2})$,
coincidindo com o termo principal da correção clássica de Stirling para uma lei
$\mathrm{Gamma}(N,1)$: após a inclinação, os termos com $b_j\ge1$ de $sX=\sum_jb_jU_j$ se comportam
assintoticamente como variáveis $\mathrm{Exp}(1)$ independentes, de modo que $sX$ sob sua inclinação
é assintoticamente distribuída como $\mathrm{Gamma}(N,1)$, e $\rho$ entra apenas através de
deslocamentos $O(1)$ dos cumulantes, em ordem $N^{-2}$. Não precisamos dessa taxa mais fina e não
provamos o termo do resto aqui; o $o(1)$ qualitativo do Teorema~\ref{thm:formulaA} é tudo que a
Seção~\ref{sec:proofs} usa.
\end{remark}

\section{O lema do envelope}
\label{sec:envelope}

O Teorema~\ref{thm:formulaA} compara $\varphi$ a $\varphi_{\mathrm{sp}}$, uma quantidade construída
a partir da transformada completa $K=Q+H+\Delta$. Resta comparar $\varphi_{\mathrm{sp}}$ a um
sistema auxiliar suave construído apenas a partir de $Q$, de modo que a parte periódica $H$ possa
ser extraída explicitamente. Escreva $t=e^{-\tau}$, $g_\tau(w) = L(w)+e^{w-\tau}$, $g_{0,\tau}(w) =
Q(w)+e^{w-\tau}$. Os sistemas verdadeiro e suave têm, cada um, seu próprio análogo de $-K'$:
escreva $B_{\mathrm{tr}}(w) := e^wt(e^w) = -L'(w)$ e $B_{\mathrm{sm}}(w) := -Q'(w) = w/c-a$.

\begin{lemma}[Selas verdadeira e suave]
\label{lem:saddles}
Sejam $\Phi(w) = w-\log B_{\mathrm{tr}}(w)$ e $\Phi_0(w) = w-\log B_{\mathrm{sm}}(w)$, esta última
definida para $w>ca$ onde $B_{\mathrm{sm}}(w)>0$, e defina
$\tau_0 := 1+ca+\log c = 0.9502067913\ldots$, de modo que $\tau_0 > \log 2$.

Para $\tau>\log 2$, $g_\tau$ tem um único minimizador global $w^*(\tau)$, no qual $g_\tau''(w^*) =
V(r^*)>0$, $r^*=e^{w^*}$; ele é caracterizado por $\tau = \Phi(w^*)$, e $\Phi$ é uma bijeção
estritamente crescente $\R\to(\log 2,\infty)$. Para $\tau\le\log 2$ não há ponto crítico algum. Para
$\tau>\tau_0$, $g_{0,\tau}$ tem um único minimizador local $w_0(\tau)$, caracterizado por
$\tau=\Phi_0(w_0)$ junto com $B_0 := B_{\mathrm{sm}}(w_0) = e^{w_0-\tau} > 1/c$; $\Phi_0$ restrita a
$(ca+1,\infty)$ é uma bijeção estritamente crescente sobre $(\tau_0,\infty)$, e $B_0$ cresce de
$1/c$ a $\infty$ ao longo dela. Para $\tau\le\tau_0$ não há tal minimizador.

Uma vez que $B_0\ge5$ (equivalentemente $\tau\ge3.73978\ldots$), $|w^*-w_0| \le 0.00652/B_0$.
Também $w_0(\tau+c)-w_0(\tau) = c+O(1/\tau)$ quando $\tau\to\infty$.
\end{lemma}

\begin{proof}
Substituir $s=e^w$ e $t=e^{-\tau}$ transforma $g_\tau$ em $K(s)+st$, estritamente convexa em $s>0$
porque $K''>0$, com derivada $K'(s)+t$ se anulando exatamente onde $t = t(s)$. Como $t(\cdot)$
decresce estritamente de $1/2$ a $0$, isso acontece para exatamente um $s$ quando $t\in(0,1/2)$,
isto é, quando $\tau>\log2$, e para nenhum $s$ quando $t\ge1/2$. Como $w\mapsto e^w$ é uma bijeção
$\R\to(0,\infty)$, $w^* = \log s$ é o único minimizador global de $g_\tau$, e $\Phi$ inverte
$\tau\mapsto w^*$. Derivando duas vezes, $g_\tau''(w) = e^w\bigl(K'(e^w)+t\bigr)+e^{2w}K''(e^w)$,
cujo primeiro termo se anula em $w^*$, deixando $g_\tau''(w^*) = V(r^*)>0$.

Para o sistema suave, $g_{0,\tau}'(w) = e^{w-\tau}-B_{\mathrm{sm}}(w)$ e $g_{0,\tau}''(w) =
e^{w-\tau}-1/c$, de modo que um ponto crítico é um minimizador local exatamente quando $B_0 > 1/c$.
Também $\Phi_0'(w) = 1-1/(c\,B_{\mathrm{sm}}(w))$, positiva exatamente onde
$B_{\mathrm{sm}}(w)>1/c$, isto é, em $w>ca+1$; ali $\Phi_0$ cresce de $\Phi_0(ca+1) = ca+1+\log c =
\tau_0$ a $\infty$. Isso dá a existência, unicidade e intervalo enunciados.

Agora tome $B_0\ge5$. A decomposição $L = Q+H+\Delta$ avalia o gradiente verdadeiro exatamente na
sela suave: $Q'(w_0) = -B_0 = -e^{w_0-\tau}$, de modo que os termos principais se cancelam e
\begin{equation}
\label{eq:gradatw0}
g_\tau'(w_0) = L'(w_0)+e^{w_0-\tau} = H'(w_0)+\Delta'(w_0), \qquad
|g_\tau'(w_0)| \le \varepsilon_1 := \sup|H'|+\sup_{[w_0-1,w_0+1]}|\Delta'| .
\end{equation}
Em $[w_0-1,w_0+1]$ escreva $g_\tau''(w) = e^{w-\tau}-1/c+H''(w)+\Delta''(w)$. Aqui $e^{w-\tau} =
B_0e^{w-w_0}$ varia sobre $[B_0/e,\,B_0e]$, de modo que $g_\tau''$ não admite cota independente de
$\tau$; ela é de ordem exata $B_0$,
\begin{equation}
\label{eq:g2scale}
\frac{B_0}{e}-\eta_0 \;\le\; g_\tau''(w) \;\le\; B_0e, \qquad
\eta_0 := \frac1c+\sup|H''|+\sup_{[w_0-1,w_0+1]}|\Delta''|,
\end{equation}
a cota superior porque $-1/c+\sup|H''|+\sup|\Delta''|<0$. Como $w_0 = c(B_0+a)$, $B_0\ge5$ força
$w_0\ge5.3492$ e portanto $w\ge4.3492$ no intervalo, onde $|\Delta'|,|\Delta''|\le10^{-60}$: com
$b:=2e^w\ge2e^{4.3492}>154$ e $x_k:=3^kb$, derivar $\Delta(w)=-\sum_{k\ge0}\log(1-e^{-x_k})$ termo
a termo (justificado como na prova do Teorema~\ref{thm:H}) dá $\Delta'(w) =
-\sum_{k\ge0}x_k/(e^{x_k}-1)$ e, aplicando $|d/dx[x/(e^x-1)]|\le8xe^{-x}$ para $x\ge1$ (elementar,
já que $e^x-1\ge e^x/2$ ali) junto com $dx_k/dw=x_k$, $\Delta''(w) = -\sum_{k\ge0}x_k\cdot
\frac{d}{dx}\Bigl[\frac{x}{e^x-1}\Bigr]_{x=x_k}$, de modo que
$|\Delta''(w)|\le\sum_{k\ge0}8x_k^2e^{-x_k}$; limitando ambas as somas pelo termo dominante $k=0$
via $x_k-b\ge2kb$ (de $3^k-1\ge2k$) dá $|\Delta'(w)|\le4be^{-b}$ e $|\Delta''(w)|\le16b^2e^{-b}$,
ambos bem abaixo de $10^{-60}$ em $b>154$. Com \eqref{eq:Hbounds} isso torna
$\varepsilon_1 \le 0.0011978$ e $\eta_0 \le 0.9170912$, donde $2e\eta_0 \le 4.9859 \le 5 \le B_0$,
isto é, $\eta_0 \le B_0/(2e)$, e \eqref{eq:g2scale} melhora para $g_\tau''\ge B_0/(2e)>0$ ao longo
de todo o intervalo.

Ponha $h_0 := 2e\varepsilon_1/B_0 \le 2e\cdot0.0011978/5 < 0.0014$. Para $h\in(h_0,1]$ o teorema do
valor médio aplicado a $[w_0-1,w_0+1]$ dá $g_\tau'(w_0+h) \ge g_\tau'(w_0)+hB_0/(2e) \ge
-\varepsilon_1 + hB_0/(2e) > 0$ e, simetricamente, $g_\tau'(w_0-h) \le \varepsilon_1-hB_0/(2e) < 0$.
Como $w^*$ é o único zero de $g_\tau'$, ele está em $(w_0-h,w_0+h)$ para todo $h$ desse tipo, de
modo que $|w^*-w_0| \le h_0 = 2e\varepsilon_1/B_0 \le 0.00652/B_0$.

Para a deriva de fase, $w_0=c(B_0+a)$ e $\tau=w_0-\log B_0$ dão $dw_0/d\tau = (1-1/(cB_0))^{-1} =
1+O(1/B_0)$ ao longo da curva $\tau\mapsto w_0(\tau)$; como $B_0\to\infty$ quando $\tau\to\infty$
(via $B_0\sim\tau/c$), integrar sobre um período dá $w_0(\tau+c)-w_0(\tau) =
\int_\tau^{\tau+c}(1+O(1/B_0(u)))\,du = c+O(1/\tau)$.
\end{proof}

\begin{proposition}[Estimativa do envelope]
\label{prop:envelope}
Com $S(\tau) := \log\varphi_{\mathrm{sp}}(t) = g_\tau(w^*)+w^*-\tfrac12\log(2\pi V(r^*))$ e
$P(\tau) := Q(w_0)+B_0+w_0-\tfrac12\log\bigl(2\pi(B_0-\tfrac1c)\bigr)$,
\[
|S(\tau) - P(\tau) - H(w_0(\tau))| = O(1/\tau) \qquad (\tau\to\infty),
\]
uniformemente na fase de $\tau$ módulo $c$.
\end{proposition}

\begin{proof}
Comece pela identidade. Como $g_{0,\tau}(w_0) = Q(w_0)+e^{w_0-\tau} = Q(w_0)+B_0$, a definição de
$P$ se lê $P(\tau) = g_{0,\tau}(w_0)+w_0-\tfrac12\log\bigl(2\pi(B_0-\tfrac1c)\bigr)$, de modo que
\[
S(\tau)-P(\tau) = \bigl[g_\tau(w^*)-g_{0,\tau}(w_0)\bigr]+(w^*-w_0)
-\frac12\log\frac{V(r^*)}{B_0-\tfrac1c}.
\]
O Teorema~\ref{thm:H} dá $L = Q+H+\Delta$, logo $g_\tau = g_{0,\tau}+H+\Delta$ pontualmente, e em
particular $g_\tau(w^*) = g_{0,\tau}(w^*)+H(w^*)+\Delta(w^*)$. Substituindo e subtraindo $H(w_0)$,
\begin{equation}
\label{eq:envfive}
\begin{split}
S(\tau)-P(\tau)-H(w_0) = {}& \underbrace{\bigl[g_{0,\tau}(w^*)-g_{0,\tau}(w_0)\bigr]}_{T_1}
+\underbrace{\bigl[H(w^*)-H(w_0)\bigr]}_{T_2}
+\underbrace{(w^*-w_0)}_{T_3} \\[2pt]
&{}\underbrace{-\;\frac12\log\frac{V(r^*)}{B_0-\tfrac1c}}_{T_4}
\;+\;\underbrace{\Delta(w^*)}_{T_5}.
\end{split}
\end{equation}
Isto é exato para todo $\tau$ com $B_0\ge5$. Os cinco termos são limitados um a um.

$T_1$ é o gap de envelope do sistema suave. Taylor em torno de $w_0$, onde $g_{0,\tau}'$ se anula,
com $g_{0,\tau}''(w) = e^{w-\tau}-1/c \le B_0e$ em $[w_0-1,w_0+1]$ e $|w^*-w_0| \le 0.00652/B_0$
pelo Lema~\ref{lem:saddles}, dá $0\le T_1 \le \tfrac12 B_0e\,(0.00652/B_0)^2 < 5.78\times10^{-5}/B_0$.

$T_2$ é a diferença da parte periódica: $|T_2| \le \sup|H'|\cdot|w^*-w_0| < 7.82\times10^{-6}/B_0$
por \eqref{eq:Hbounds}.

$T_3$ é o próprio deslocamento da localização de sela, $|T_3| \le 0.00652/B_0$, novamente pelo
Lema~\ref{lem:saddles}.

$T_4$ compara as duas normalizações. Por $V(r^*) = g_\tau''(w^*)$ e a expressão para $g_\tau''$
usada no Lema~\ref{lem:saddles},
\[
V(r^*)-\Bigl(B_0-\frac1c\Bigr) = B_0\bigl(e^{w^*-w_0}-1\bigr)+H''(w^*)+\Delta''(w^*),
\]
cujo módulo é no máximo $0.00652e^{0.00652/5}+\sup|H''|+10^{-60} < 0.0135$, já que
$B_0|e^{w^*-w_0}-1| \le B_0|w^*-w_0|e^{|w^*-w_0|}$. Dividindo por $B_0-1/c$ e tomando
$\tfrac12\log(1+\cdot)$ dá $|T_4| \le 0.0068/(B_0-1/c)$.

$T_5$ é o resto do Teorema~\ref{thm:H}, avaliado na sela verdadeira em vez de descartado. Aqui
$B_0\ge5$ força $w_0\ge5.3492$ e $w^*\ge w_0-0.0014$, de modo que $e^{w^*}\ge210.1$ e
$|T_5| \le \sum_{k\ge0}2e^{-2\cdot3^k\cdot210.1} < 10^{-182}$.

Cada cota é $O(1/B_0)$ (para $T_4$, via $B_0-1/c\ge B_0-1$, comparável a $B_0$ uma vez que
$B_0\ge5$). Finalmente $\tau = w_0-\log B_0$ e $w_0 = c(B_0+a)$ dão $\tau = cB_0+ca-\log B_0$, de
modo que $B_0/\tau \to 1/c$, e $O(1/B_0) = O(1/\tau)$ com constantes livres da fase de $\tau$.
\end{proof}

\section{A identidade de Berg--Kr\"uppel}
\label{sec:bk}

\begin{proposition}
\label{prop:bkidentity}
Escreva $f(p) = \alpha\log p-\beta\log^2 p$, de modo que a transformada $G_0(p) =
p^\alpha\exp(-\beta\log^2 p)$ de Berg e Kr\"uppel, de sua equação~(9.3), é $e^{f(p)}$ e sua
equação~(9.5) se lê
\[
g_0(t) = \frac{1}{2\pi i}\int_{\sigma-i\infty}^{\sigma+i\infty}\exp[pt+f(p)]\,dp, \qquad \sigma>0.
\]
Tome $\alpha$, $\beta$ de sua equação~(9.4) em $a=3$, $\lambda=2/3$. Então $f(p) = Q(\log p)$
identicamente; a sela exata dessa integral, $t+f'(x)=0$ com $t = e^{-\tau}$, é $x^\sharp =
e^{w_0(\tau)}$; e
\[
P(\tau) \;=\; x^\sharp t + f(x^\sharp) - \tfrac12\log\bigl(2\pi f''(x^\sharp)\bigr),
\]
a expressão de ponto de sela invocada na prova de sua Proposição~9.1, avaliada nessa sela exata. Na
variável $Y := cB_0$, com $r = -ca-\log c$, a equação de sela é $Y-\log Y = \tau+r$.
\end{proposition}

\begin{proof}
A equação truncada deles~(9.2) é $\lambda \psi'(t) = a\,\psi(at)$ para $a>1$, $\lambda>0$;
transformando por Laplace (o termo de contorno $\lambda\psi(0)$ de $\mathcal L[\psi'](p)=pG(p)-\psi(0)$
se anula, já que as soluções de interesse são suportadas em $[0,\infty)$ com $\psi(0)=0$) dá $\lambda
p\,G(p) = G(p/a)$, cuja solução geral~(9.3) é $G_0$ vezes uma função $1$-periódica arbitrária de
$\log p/\log a$, com
\[
\alpha = -\frac12-\frac{\log\lambda}{\log a}, \qquad \beta = \frac{1}{2\log a}
\]
sua equação~(9.4). A Seção~\ref{sec:prelim} colocou $\varphi$ nessa família em $a=3$, $\lambda=2/3$:
a equação~(9.1) não-truncada deles se lê $\tfrac23\varphi'(t) = 3\bigl(\varphi(3t) -
\varphi(3t-2)\bigr)$, que em $(0,\tfrac23)$ é $\varphi'(t) = \tfrac92\varphi(3t)$. Nesses valores
\[
\alpha = -\frac12-\frac{\log(2/3)}{\log3} = \frac12-\frac{\log2}{\log3} = a, \qquad
\beta = \frac{1}{2\log3} = \frac{1}{2c},
\]
onde o $a$ à direita é o coeficiente linear de $Q$ deste artigo; o $a$ deles, o fator de dilatação
$3$, é uma quantidade diferente e não é escrito novamente abaixo. Logo, escrevendo $y := \log p$,
\[
f(p) = \alpha y-\beta y^2 = a y - \frac{y^2}{2c} = Q(y),
\]
a primeira das quatro afirmações da proposição. Esta é uma identidade em $y$, não apenas uma
coincidência de valor: a regra da cadeia ($p\,d/dp = d/dy$) aplicada às primeira e segunda
derivadas de $f$ confirma isso também no nível de $Q'$ e $Q''$,
\[
-p\,f'(p) = -Q'(y) = B_{\mathrm{sm}}(y) = 2\beta y-\alpha, \qquad
p^2f''(p) = Q''(y)-Q'(y) = B_{\mathrm{sm}}(y)-\frac1c = 2\beta y-2\beta-\alpha .
\]
A própria fórmula de Berg e Kr\"uppel para $f''(p)$, exibida na p.~179 de \cite{bergkruppel1998},
na prova de sua Proposição~9.1 (a exibição não numerada logo após ``In view of''), se lê
$p^2f''(p)=2\beta\ln p+2\beta-\alpha$, com $+2\beta$ onde a diferenciação direta de $f$ dá $-2\beta$:
$f''(p) = \tfrac{1}{p^2}\bigl(2\beta\ln p - 2\beta - \alpha\bigr)$, já que
$f'(p)=\tfrac1p(\alpha-2\beta\ln p)$ se deriva para $-\tfrac{2\beta}{p^2}-\tfrac{\alpha-2\beta\ln
p}{p^2}$. A Proposição~9.1 deles não é afetada, sendo o termo de ordem menor ali do que
$2\beta\log p$, mas a expressão exata usada aqui é afetada: a discrepância $\pm2\beta$ desloca
$(x^\sharp)^2f''(x^\sharp) = B_0-1/c$ por $2/c$, portanto desloca $-\tfrac12\log\bigl(2\pi
(B_0-1/c)\bigr)$ em $P(\tau)$ abaixo por $O(1/B_0)$, que se anula quando $B_0\to\infty$.

A condição de sela $t+f'(x)=0$ em $x=e^y$ agora se lê $e^{-\tau} = B_{\mathrm{sm}}(y)e^{-y}$, isto
é, $\tau = y-\log B_{\mathrm{sm}}(y) = \Phi_0(y)$, cuja única solução no ramo $B_{\mathrm{sm}}>1/c$
é $y = w_0(\tau)$ pelo Lema~\ref{lem:saddles}. Assim $x^\sharp = e^{w_0}$, e ali $x^\sharp t =
e^{w_0-\tau} = B_0$, $f(x^\sharp) = Q(w_0)$, $(x^\sharp)^2f''(x^\sharp) = B_0 - 1/c>0$. Substituindo,
\[
x^\sharp t+f(x^\sharp)-\tfrac12\log\bigl(2\pi f''(x^\sharp)\bigr)
= B_0+Q(w_0)-\tfrac12\log\frac{2\pi(B_0-1/c)}{(x^\sharp)^2}
\]
\[
= Q(w_0)+B_0+w_0-\tfrac12\log\Bigl(2\pi\bigl(B_0-\tfrac1c\bigr)\Bigr),
\]
que é $P(\tau)$. Finalmente $w_0 = c(B_0+a)$ e $\tau = w_0-\log B_0$ dão $\tau = cB_0+ca-\log
B_0 = Y-\log Y+ca+\log c = Y-\log Y-r$.
\end{proof}

Revertendo em série $Y-\log Y=\tau+r$ até segunda ordem: escrevendo $W:=\tau+r$, uma primeira
substituição $Y=W+\log W+\delta$ em $Y-\log Y=W$ dá, já que $\log(1+u)=u+O(u^2)$ com
$u=(\log W+\delta)/W$, $\delta = (\log W+\delta)/W + O(\log^2W/W^2)$, de modo que
$\delta=\log(W)/W+O(\log^2W/W^2)$ (o termo único $Y=W+\log W$ sozinho só é preciso a
$O(\log W/W)$, uma ordem grosseira demais; este segundo termo é necessário). Logo
$Y=W+\log W+\log(W)/W+O(\log^2W/W^2)$, e $\log W=\log\tau+r/\tau+O(1/\tau^2)$ então dá
$Y = \tau + \log\tau + r + (\log\tau+r)/\tau + O(\tau^{-2}\log^2\tau)$. Como $w_0=Y+ca$ e
$Q(w_0)=-Y^2/(2c)+ca^2/2$ (substituindo $w_0=Y+ca$ na definição de $Q$ e simplificando), o termo
principal de $P$ em $Y$ é $-Y^2/(2c)$, de modo que $\partial P/\partial Y=-Y/c+O(1)$ transforma o
erro $O(\tau^{-2}\log^2\tau)$ em $Y$ acima no erro $O(\tau^{-1}\log^2\tau)$ em $P$ abaixo.
Substituindo em $P$ dá
\[
P(\tau) = -\frac{(\tau+\log\tau)^2}{2c} + \Bigl(1+\tfrac1c-\tfrac{r}{c}\Bigr)\tau +
\Bigl(\tfrac12-\tfrac{r}{c}\Bigr)\log\tau + C_P + O\!\left(\frac{\log^2\tau}{\tau}\right),
\]
\[
C_P = -\frac{r^2}{2c}+r+ca+\frac{ca^2}{2}-\frac{\log2\pi}{2}+\frac{\log c}{2} =
-0.9576743183133982760669122497855855\ldots.
\]
As próprias equações de Berg e Kr\"uppel, seguintes à~(9.6), dão, na notação deles,
\[
\delta_{\mathrm{BK}} = \tfrac12+\alpha-2\beta\log(2\beta), \qquad
\gamma = -2\beta-\delta_{\mathrm{BK}}-\tfrac12, \qquad
\varepsilon = \tfrac12+\alpha-\beta\log(2\beta),
\]
com $\alpha=a$, $\beta=1/(2c)$ como identificado acima. Substituindo e simplificando (usando
$r=-ca-\log c$) dá, coincidindo os coeficientes de $\tau$, $\log\tau$ e a constante de $P$ termo a
termo,
\[
\gamma = -\Bigl(1+\tfrac1c-\tfrac rc\Bigr) = -1.8649154949\ldots, \qquad
\delta_{\mathrm{BK}} = \tfrac12-\tfrac rc = 0.4546762683\ldots,
\]
\[
\varepsilon = a+\tfrac12+\frac{\log c}{2c} = 0.4118732574\ldots,
\]
de modo que, em particular, $C_P=\varepsilon\log(2\beta)-\tfrac12\log(2\pi)$ vale identicamente
(verificado substituindo a expressão de $\varepsilon$ diretamente na de $C_P$), dando
\[
C_P = \varepsilon\log(2\beta) - \tfrac12\log(2\pi) = \log\frac{(2\beta)^\varepsilon}{\sqrt{2\pi}}.
\]
Logo $P(\tau) - \log\varphi_{0,\mathrm{bare}}(e^{-\tau}) \to C_P$ e, usando a normalização com o
prefator incluído, $P(\tau) - \log\varphi_0(e^{-\tau}) \to 0$ com taxa $O(\tau^{-1}\log^2\tau)$: os
coeficientes de $\tau$, $\log\tau$, e o termo constante de $P$ reproduzem $\gamma$,
$\delta_{\mathrm{BK}}$, $\varepsilon$ exatamente, de modo que essa convergência é consequência
direta da identidade da Proposição~\ref{prop:bkidentity} junto com a reversão de série acima, não um
apelo à própria análise assintótica de Berg e Kr\"uppel para $\varphi_0$.

\begin{proof}[Prova do Corolário~\ref{cor:bk}]
Pelas Proposições~\ref{prop:envelope} e~\ref{prop:bkidentity}, $S(\tau) = \log\varphi_0(e^{-\tau}) +
H(w_0(\tau)) + O(\tau^{-1}\log^2\tau)$, isto é, $\varphi_{\mathrm{sp}} = \varphi_0\cdot
e^{H(w_0(\tau))}(1+o(1))$; pelo Teorema~\ref{thm:formulaA}, $\varphi=\varphi_{\mathrm{sp}}(1+o(1))$
também, de modo que $\varphi = \varphi_0\cdot e^{H(w_0(\tau))}(1+o(1))$: o fator periódico que
multiplica a solução assintótica de Berg e Kr\"uppel, na equação~(9.7) para esta autofunção, é
$e^H$, coincidindo com a fórmula de produto que sua Proposição~9.3 fornece para ele, como observado
após o Teorema~\ref{thm:H}. A série de Fourier de $H$ da Proposição~\ref{prop:Hfourier} é a forma
fechada nova; o certificado da Proposição~\ref{prop:Hcert} é o que nem a fórmula de produto deles
nem sua observação final (de que eles apenas \emph{esperam} que $\varphi$ seja limitada e limitada
longe de zero relativamente a $\varphi_0$) resolvem.
\end{proof}

\section{Provas dos teoremas principais}
\label{sec:proofs}

Ao longo desta seção, o parâmetro $N$ do Teorema~\ref{thm:formulaA} e o parâmetro $\tau$ das
Seções~\ref{sec:envelope}--\ref{sec:bk} são combinados livremente num único $o(1)$: já que
$s=e^{w^*}$ e, pelo Lema~\ref{lem:saddles}, $w^*=w_0(\tau)+O(1/B_0)$ com $w_0(\tau)=c(B_0+a)$,
$N=\lfloor\log_3(2s)\rfloor \sim w^*/c \sim \tau/c$ quando $\tau\to\infty$, de modo que
$N\to\infty$ exatamente quando $\tau\to\infty$. Em $\widetilde{A}_\delta$ abaixo, $\tau_l\to\infty$
uniformemente, a uma taxa que depende apenas de $l$ e $\delta$, não da sequência particular
escolhida.

\begin{proof}[Prova do Teorema~\ref{thm:conj3}]
A classe $\widetilde{A}_\delta$ de Wirsching (suas equações~(1.5) e~(3.2)) consiste de sequências
$(x_l)$ com $\lfloor 3^l x_l\rfloor = k_l$ satisfazendo $|l-k_l|\le\delta\sqrt l$ para todo $l$; já
que $|3^lx_l-k_l|<1$ também, $|3^lx_l-l|<1+\delta\sqrt l$, logo $\lambda_l := 3^lx_l/l \to 1$
uniformemente a taxa $O(1/l)+O(\delta/\sqrt l)$ na classe. Escrevendo $z_l :=
\lambda_l\, l\, 3^{-l} = x_l$ (o próprio $z_l$ de Wirsching, do enunciado da conjectura, e seu
$x_l$, da definição da classe, coincidem, por definição de $\lambda_l$), $\tau_l = -\log z_l = lc -
\log l - \log\lambda_l$, a localização de sela satisfaz $w_0(\tau_l) = lc - \log\lambda_l + O(1/l)$:
a reversão de $Y-\log Y=\tau+r$ da Seção~\ref{sec:bk} dá
$w_0(\tau) = \tau+\log(\tau/c)+(\log\tau+r)/\tau+O(\tau^{-2}\log^2\tau)$, e o truncamento ingênuo de
um único termo $w_0(\tau)\approx\tau+\log(\tau/c)$ sozinho só é preciso a $O(\log\tau/\tau)$, uma
ordem grosseira demais; substituindo $\tau=\tau_l$ na fórmula completa de dois termos, as partes
$O(\log l/l)$ ao expandir $\log(\tau_l/c)$ contra $\tau_l = lc-\log l-\log\lambda_l$ se cancelam
exatamente contra o próprio termo $(\log\tau_l+r)/\tau_l$ da reversão, deixando o $O(1/l)$
enunciado, o mesmo mecanismo de cancelamento da reversão de série acima. Assim
$\operatorname{dist}(w_0(\tau_l), c\Z) \to 0$ uniformemente na classe, a taxa $O(1/\sqrt l)$:
escrevendo $w_0(\tau_l) = lc + \varepsilon_l$ com $\varepsilon_l = -\log\lambda_l + O(1/l) \to 0$, o
múltiplo mais próximo de $c$ é o próprio $lc$ uma vez que $|\varepsilon_l|<c/2$, de modo que
$H(w_0(\tau_l)) \to H(0)$ por continuidade, independentemente do sinal de $\varepsilon_l$. Pelo
Teorema~\ref{thm:formulaA}, Proposição~\ref{prop:envelope}, Proposição~\ref{prop:bkidentity} (que
dá $P(\tau)-\log\varphi_0(e^{-\tau})\to0$, Seção~\ref{sec:bk}), e continuidade de $H$ (de fato
Lipschitz, por \eqref{eq:Hbounds}),
\[
\log\varphi(z_l) - \log\varphi_0(z_l) = H(w_0(\tau_l)) + o(1) \to H(0)
\]
uniformemente em $\widetilde{A}_\delta$, isto é, $\varphi(z_l)/\varphi_0(z_l) \to e^{H(0)}$
uniformemente. Isso é exatamente a Conjectura 3, com a constante dada explicitamente.
\end{proof}

\begin{remark}
A prova acima usa apenas $\lambda_l\to1$, equivalentemente $|l-k_l|=o(l)$; a cota específica
$O(\sqrt l)$ de Wirsching em $\widetilde{A}_\delta$ não desempenha papel algum além de forçar essa
condição mais fraca. A mesma conclusão, e o mesmo limite $e^{H(0)}$, valem para a classe mais ampla
de sequências satisfazendo apenas $|l-k_l|=o(l)$.
\end{remark}

\begin{proof}[Prova do Teorema~\ref{thm:unrestricted}]
Como $t=e^{-\tau}$ percorre $(0,1/2)$ sem restrição, $\tau\to\infty$. Pelo Lema~\ref{lem:saddles},
$w_0=\Phi_0^{-1}$ é contínua, estritamente crescente e ilimitada em $(\tau_0,\infty)$, de modo que
$w_0(\tau)\bmod c$ varre todo valor em $[0,c)$ infinitas vezes quando $\tau\to\infty$ (uma função
contínua, estritamente crescente e ilimitada atinge toda classe residual módulo $c$ infinitas
vezes), de modo que pela mesma combinação de resultados usada acima,
$\log\varphi(t)-\log\varphi_0(t) = H(w_0(\tau))+o(1)$ assume valores arbitrariamente próximos tanto
de $\sup H$ quanto de $\inf H$ infinitas vezes. Pela Proposição~\ref{prop:Hcert},
$\operatorname{osc}(H) > 0$ rigorosamente, de modo que $\varphi(t)/\varphi_0(t)$ não converge, e,
escrevendo $\Psi(t) := \varphi(t)/\varphi_0(t)$,
\[
\frac{\limsup_{t\to0^+}\Psi(t)}{\liminf_{t\to0^+}\Psi(t)} \ge e^{\operatorname{osc}(H)} \ge 1.0004188:
\]
o limsup da razão excede seu liminf por um fator de pelo menos $1.0004188$, isto é, por pelo menos
$0.0418\%$.
\end{proof}

\begin{proposition}[A exigência real de Wirsching]
\label{prop:wirschingreq}
Para $(x_l)\in\widetilde A_\delta$, escreva $x_l^+ := x_l + 3^{-(l+1)}$ (sua própria definição,
prosa não numerada na p.~16 de \cite{wirsching2003}, imediatamente antes de sua equação~(7.14);
como $x_l\sim l\cdot3^{-l}$, $x_l^+$ é assintótico a $x_l$ e desempenha o mesmo papel que o $z_l$
deste artigo). O próprio uso de Wirsching da Conjectura 3, via $\varphi'(x)=\tfrac92\varphi(3x)$ em
$(0,\tfrac23)$ e seu próprio cálculo, sua equação~(7.13), de que
$\lim_l (2/3^{l+1})\varphi_0'(x_l)/\varphi_0(x_l) = 2/3$ uniformemente na classe, se reduz a
\[
\limsup_l \Lambda_l < \tfrac32, \qquad
\Lambda_l := \frac{\varphi(3x_l^+)/\varphi_0(3x_l^+)}{\varphi(x_l^+)/\varphi_0(x_l^+)}.
\]
Isso vale com ampla margem: $\Lambda_l \to 1$.
\end{proposition}

\begin{proof}
Primeiro, a redução em si. A exigência real de Wirsching é sua desigualdade~(7.5),
\[
\limsup_l \frac{2}{3^{l+1}}\cdot\frac{\varphi'(x_l^+)}{\varphi(x_l^+)} < 1;
\]
por $\varphi'=\tfrac92\varphi(3\cdot)$ isso é $\limsup_l (9/3^{l+1})\varphi(3x_l^+)/\varphi(x_l^+) <
1$. Pela definição de $\Lambda_l$,
$\varphi(3x_l^+)/\varphi(x_l^+) = \Lambda_l\cdot\varphi_0(3x_l^+)/\varphi_0(x_l^+)$, e o $\varphi_0$
de Berg e Kr\"uppel é uma solução assintótica da mesma equação truncada (a equação~(7.12) deles, na
notação de Wirsching: $\varphi_0'(t)/\varphi_0(3t)\to9/2$ quando $t\to0$), de modo que
$\varphi_0(3x_l^+)/\varphi_0(x_l^+) = \tfrac29\varphi_0'(x_l^+)/\varphi_0(x_l^+)\,(1+o(1))$.
Para ver que (7.13) se transfere de $x_l$ para $x_l^+$ com uma correção $o(1)$, escreva $v=\log t$;
a forma explícita de $\varphi_0$ acima dá $\varphi_0'(t)/\varphi_0(t) = B(v)/t$ com
$B(v) = \gamma+\delta_{\mathrm{BK}}/v-2\beta(v-\log(-v))(1-1/v)$, e, quando $v\to-\infty$,
$B(v)\sim-2\beta v$ enquanto $B'(v)\to-2\beta$, de modo que $(\log B)'(v)=B'(v)/B(v)\to0$ e
$(\log[\varphi_0'/\varphi_0])'(v) = (\log B)'(v)-1 \to -1$, limitado. Como
$x_l^+/x_l = 1+3^{-(l+1)}/x_l \to 1$ (já que $x_l\sim l\cdot3^{-l}$, de fato
$3^{-(l+1)}/x_l=O(1/l)$), $v(x_l^+)-v(x_l)=\log(x_l^+/x_l)=O(1/l)$, e o teorema do valor médio dá
$\log[\varphi_0'/\varphi_0(x_l^+)]-\log[\varphi_0'/\varphi_0(x_l)]=O(1/l)$, isto é,
$\varphi_0'(x_l^+)/\varphi_0(x_l^+)=\varphi_0'(x_l)/\varphi_0(x_l)\cdot(1+O(1/l))$. Substituindo e
usando sua equação~(7.13), $\lim_l(2/3^{l+1})\varphi_0'(x_l)/\varphi_0(x_l)=2/3$ uniformemente na
classe,
\[
\frac{9}{3^{l+1}}\cdot\frac{\varphi(3x_l^+)}{\varphi(x_l^+)}
= \Lambda_l\cdot\frac2{3^{l+1}}\cdot\frac{\varphi_0'(x_l^+)}{\varphi_0(x_l^+)}\,(1+o(1))
= \frac23\,\Lambda_l\cdot(1+o(1)),
\]
de modo que a exigência $\limsup_l(9/3^{l+1})\varphi(3x_l^+)/\varphi(x_l^+)<1$ é exatamente
$\limsup_l\Lambda_l<3/2$.

Para o valor de $\Lambda_l$: escreva $\tau_l := -\log x_l^+$ (localmente a esta prova apenas,
distinto do $\tau_l$ da prova do Teorema~\ref{thm:conj3}), de modo que $t=x_l^+$ corresponde a
$\tau=\tau_l$ e $t=3x_l^+$ a $\tau=\tau_l-c$. O deslocamento de fase sob $t\mapsto3t$, isto é,
$\tau\mapsto\tau-c$, é $w_0(\tau-c)-w_0(\tau)+c = O(1/\tau)$ pela equação que define o
Lema~\ref{lem:saddles}, independentemente da fase de $\tau$; como $H$ é $c$-periódica,
$H(w_0(\tau-c)) = H(w_0(\tau-c)+c) = H(w_0(\tau)+O(1/\tau))$, e $H$ Lipschitz (por
\eqref{eq:Hbounds}) transforma isso em $H(w_0(\tau-c)) = H(w_0(\tau)) + O(1/\tau)$. Pelo
Teorema~\ref{thm:formulaA}, Proposição~\ref{prop:envelope}, e Proposição~\ref{prop:bkidentity}
juntas (a mesma cadeia usada na prova do Teorema~\ref{thm:conj3}),
$\log\varphi(t)-\log\varphi_0(t) = H(w_0(\tau)) + o(1)$ uniformemente em fase, de modo que
\[
\log\Lambda_l = H(w_0(\tau_l-c)) - H(w_0(\tau_l)) + o(1) = O(1/\tau_l) + o(1) \to 0,
\]
usando que $H$ é Lipschitz (\eqref{eq:Hbounds}) e o deslocamento de fase acima é $O(1/\tau_l)$;
logo $\Lambda_l \to 1$, bem dentro de $\limsup \Lambda_l < 3/2$.
\end{proof}

\section{Discussão}
\label{sec:discussion}

Estes resultados resolvem apenas a Conjectura 3. O argumento de densidade positiva de predecessores
de Wirsching se reduz a três conjecturas; as Conjecturas 1 e 2, tratando respectivamente de uma
transferência pontual-versus-média entre os geradores $g_l(k,a)$ e sua média de Haar $3$-ádica
(substituídos pelas próprias funções Elka, definidas em \cite{wirsching1998}, Seção II.4, para
tornar a normalização possível) e de uma cota inferior positiva uniforme numa razão de iterados de
$W_3$ ao longo de $\widetilde A_\delta$, permanecem intocadas aqui, e o teorema completo não é
estabelecido por este artigo.

Wirsching enunciou a Conjectura 3, mais fraca, em vez da mais forte $\varphi\sim\kappa\varphi_0$
que ele considera primeiro. O Teorema~\ref{thm:unrestricted} mostra o porquê: o enunciado mais
forte falha por um mecanismo com amplitude certificada e positiva. Toda fase $\theta\in[0,c)$
varrida por $H$ admite alguma sequência ao longo da qual a razão converge para $e^{H(\theta)}$:
tome $\tau_l := \Phi_0(\theta+lc)$, $z_l := e^{-\tau_l}$, de modo que $w_0(\tau_l)=\theta+lc$
exatamente, e aplique o Teorema~\ref{thm:formulaA}, a Proposição~\ref{prop:envelope}, e a
Proposição~\ref{prop:bkidentity} como na prova do Teorema~\ref{thm:conj3}. O que é especial na
fase $0$ é apenas que ela é a que a própria classe $\widetilde A_\delta$ de Wirsching acaba
fixando. Essa fase está próxima de $w_{1011118}$, um dos quatro pontos onde $H$ é impresso na prova
da Proposição~\ref{prop:Hcert}: combinando esse valor com $H(0)$ dá
$H(0) - H(w_{1011118}) = 5.25\times10^{-6}$, um cálculo em ponto flutuante (não, em si, um
envolvimento certificado), contra uma oscilação certificada de pelo menos $4.18\times10^{-4}$.
(Se $w_{1011118}$ é onde $H$ atinge seu verdadeiro mínimo não é certificado, pela própria ressalva
da Proposição~\ref{prop:Hcert} sobre a busca em grade.) Não sabemos se essa quase-coincidência com
a fase $0$ é acidental.

A taxa não é estabelecida além de $o(1)$. A correção de Edgeworth mais fina da
Observação~\ref{rem:edgeworth}, e uma forma correspondentemente mais fina das comparações de
envelope e de Berg--Kr\"uppel, dariam uma taxa polinomial explícita ao longo de tudo; nenhum dos
dois teoremas precisa disso.

\section{Disponibilidade de código e dados}
\label{sec:code}

Toda afirmação certificada ou numericamente verificada neste artigo é sustentada por um script no
repositório que o acompanha, \url{https://github.com/faculdade/wirsching-conjecture3-proof},
organizado por seção com um README por script informando o que ele verifica e como executá-lo.
Isso inclui os certificados de aritmética de intervalos por trás da Proposição~\ref{prop:Hcert} e
das cotas \eqref{eq:Hbounds}, a cadeia de constantes por trás do Teorema~\ref{thm:formulaA}, e a
identidade da Proposição~\ref{prop:bkidentity}. Os scripts de aritmética de intervalos usam
aritmética de bolas de Arb via \texttt{python-flint}, com precisão de trabalho de $90$ a $100$
dígitos decimais (aproximadamente $300$ a $332$ bits, indicada exatamente no topo de cada script),
estritamente maior que os $250$ bits de precisão citados no texto na prova da
Proposição~\ref{prop:Hcert} para os valores ali impressos: os scripts arquivados são uma
reprodução independente e de maior precisão do mesmo certificado, não uma aproximação de precisão
menor. A precisão de trabalho é escolhida com margem bem além do raio final de cada certificado;
todo script desse tipo imprime suas bolas de entrada e seu envolvimento final, de modo que a saída
impressa da execução é o próprio certificado. A versão do repositório referenciada ao longo do
artigo está arquivada, independentemente de quaisquer mudanças futuras no ramo padrão do
repositório, em \href{https://doi.org/10.5281/zenodo.21854549}{doi:10.5281/zenodo.21854549}
(commit \texttt{f8248c3}).

\subsection*{Agradecimentos}
Foi usada assistência de IA para revisão textual e tradução.

\begin{thebibliography}{99}

\bibitem{wirsching2003}
G.~J. Wirsching, \emph{On the problem of positive predecessor density in $3N+1$ dynamics},
Discrete Contin. Dyn. Syst. \textbf{9} (2003), no.~3, 771--787, doi:10.3934/dcds.2003.9.771.

\bibitem{bergkruppel1998}
L.~Berg and M.~Kr\"uppel, \emph{On the solution of an integral-functional equation with a
parameter}, Z. Anal. Anwendungen \textbf{17} (1998), no.~1, 159--181.

\bibitem{wirsching1998}
G.~J. Wirsching, \emph{The Dynamical System Generated by the $3n+1$ Function}, Lecture Notes in
Mathematics, vol.~1681, Springer, Berlin, 1998.

\bibitem{rvachev1990}
V.~A. Rvachev, \emph{Compactly supported solutions of functional-differential equations and their
applications}, Russian Math. Surveys \textbf{45} (1990), no.~1, 87--120.

\bibitem{debruijn1948}
N.~G. de Bruijn, \emph{On Mahler's partition problem}, Nederl. Akad. Wetensch., Proc.
\textbf{51} (1948), 659--669 (also Indag. Math. \textbf{10}, 210--220).

\bibitem{erdosrichmond1976}
P.~Erd\H{o}s and L.~B. Richmond, \emph{Concerning periodicity in the asymptotic behaviour of
partition functions}, J. Austral. Math. Soc. Ser.~A \textbf{21} (1976), 447--456.

\bibitem{katomcleod1971}
T.~Kato and J.~B. McLeod, \emph{The functional-differential equation $y'(x) = ay(\lambda x) +
by(x)$}, Bull. Amer. Math. Soc. \textbf{77} (1971), no.~6, 891--937.

\bibitem{derfel1995}
G.~Derfel, \emph{Functional-differential and functional equations with rescaling}, in Operator
Theory: Advances and Applications, vol.~80, Birkh\"auser, Basel, 1995, 100--111.

\bibitem{titchmarsh1986}
E.~C. Titchmarsh, \emph{The Theory of the Riemann Zeta-Function}, 2nd ed., revised by D.~R.
Heath-Brown, Oxford University Press, 1986, Chapter III, \S3.1--3.3 (the non-vanishing of $\zeta$
on $\operatorname{Re}s=1$; not the quantitative zero-free region of Theorem 3.8, a stronger and
unnecessary statement for this paper's purposes).

\end{thebibliography}

\end{document}
